\section{Effective curvature and effective smoothness along the trajectory (hypothesis B2)}
\label{sec:B2}

The proof of Proposition~5.9 of \cite{bach2024learning} (convergence of SAGA, transcribed in the B3 note)
uses exactly two problem constants. The per-sample smoothness constant $R^2$ enters through the
co-coercivity inequality (equation (5.31) of the book),
\[
\frac1n\sum_{i=1}^n \norm{\fip(\theta_{t-1})-\fip(\etastar)}^2 \le R^2\,\Fp(\theta_{t-1})^\top(\theta_{t-1}-\etastar),
\]
and the strong-convexity constant $\mu$ enters through
$(\theta_{t-1}-\etastar)^\top \Fp(\theta_{t-1})\ge \mu\norm{\theta_{t-1}-\etastar}^2$.
Both inequalities are applied at the \emph{current iterate} $\theta_{t-1}$, always with the same second
point $\etastar$. Hypothesis (B2) of the project is that the slack of these two inequalities along the
trajectory actually followed by the algorithm explains a large part of the gap between the rate
$1-\min\{1/(3n),3\mu/(16R^2)\}$ and the observed contraction. This note defines the two
trajectory-dependent constants that make the two inequalities equalities, proves their basic properties,
makes them explicit for $\ell_2$-regularized logistic regression, analyses their dependence on the
direction of the error $\theta_t-\etastar$ (which can differ between gradient descent and SAGA, and which is
constrained by an exact invariance property: a trajectory started in the span of the features never leaves it,
Lemma~\ref{lem:B2-span}), and finally
explains how they enter the SAGA bound and how to measure the fraction of the theory--practice gap they
explain.

\subsection{Setting and definitions}

\begin{assumption}[convex $R^2$-smooth $f_i$, $\mu$-strongly convex $F$]\label{ass:B2}
\begin{enumerate}[label=(\roman*)]
\item Each $f_i:\R^d\to\R$ is convex, differentiable and $R^2$-smooth, i.e.\ $\fip$ is $R^2$-Lipschitz-continuous.
\item $F=\frac1n\sum_{i=1}^n f_i$ is $\mu$-strongly convex with $\mu>0$; its (unique) minimizer is
$\etastar$, so that $\Fp(\etastar)=0$ (Proposition~5.2 of \cite{bach2024learning}). We denote by $L\le R^2$
the smoothness constant of $F$ ($\Fp$ is $L$-Lipschitz-continuous; $L\le R^2$ because an average of
$R^2$-Lipschitz maps is $R^2$-Lipschitz), and by $\kappa=R^2/\mu$ the condition number used in the book.
\item Whenever second derivatives appear (Lemma~\ref{lem:B2-integral}, Proposition~\ref{prop:B2-limits},
Section~\ref{sec:B2-logistic} and after), each $f_i$ is twice continuously differentiable.
\end{enumerate}
\end{assumption}

Throughout, for $\theta\neq\etastar$ we write
\[
v=\theta-\etastar,\qquad u=v/\norm{v},
\]
and we use repeatedly that $\Fp(\theta)=\Fp(\theta)-\Fp(\etastar)=\frac1n\sum_i[\fip(\theta)-\fip(\etastar)]$.

\begin{definition}[effective constants $\muEff$ and $\LEff$]\label{def:B2-eff}
For $\theta\neq\etastar$, define
\begin{equation}\label{eq:B2-def-mueff}
\muEff(\theta)\;=\;\frac{(\theta-\etastar)^\top \Fp(\theta)}{\norm{\theta-\etastar}^2},
\end{equation}
\begin{equation}\label{eq:B2-def-Leff}
\LEff(\theta)\;=\;\frac{\frac1n\sum_{i=1}^n\norm{\fip(\theta)-\fip(\etastar)}^2}{(\theta-\etastar)^\top \Fp(\theta)}.
\end{equation}
Along a trajectory $(\theta_t)_{t\ge0}$ we abbreviate $\muEff(t)=\muEff(\theta_t)$ and $\LEff(t)=\LEff(\theta_t)$.
\end{definition}

Lemma~\ref{lem:B2-mueff} below shows that $(\theta-\etastar)^\top\Fp(\theta)\ge\mu\norm{\theta-\etastar}^2>0$,
so both quantities are well defined and positive. By construction the two inequalities quoted above
become \emph{identities} when $\mu$ and $R^2$ are replaced by $\muEff(\theta)$ and $\LEff(\theta)$:
\[
(\theta-\etastar)^\top\Fp(\theta)=\muEff(\theta)\norm{\theta-\etastar}^2,\qquad
\frac1n\sum_{i=1}^n\norm{\fip(\theta)-\fip(\etastar)}^2=\LEff(\theta)\,(\theta-\etastar)^\top\Fp(\theta).
\]
We call $\muEff(\theta)$ the \emph{effective curvature} and $\LEff(\theta)$ the \emph{effective
co-coercivity (or effective smoothness) constant} at $\theta$.

\subsection{General properties}\label{sec:B2-general}

We first record the two-point inequalities that we use, with proofs, since the whole note rests on them.

\begin{lemma}[two-point inequalities]\label{lem:B2-twopoint}
Let $f:\R^d\to\R$ be convex and differentiable.
\begin{enumerate}[label=(\alph*)]
\item (Monotonicity) For all $\theta,\eta$: $(f'(\theta)-f'(\eta))^\top(\theta-\eta)\ge0$.
\item (Strong monotonicity) If $f$ is $\mu$-strongly convex, then for all $\theta,\eta$:
$(f'(\theta)-f'(\eta))^\top(\theta-\eta)\ge\mu\norm{\theta-\eta}^2$.
\item (Co-coercivity, Proposition~5.4 of \cite{bach2024learning}) If $f$ is $L_f$-smooth, then for all $\theta,\eta$:
$\norm{f'(\theta)-f'(\eta)}^2\le L_f\,(f'(\theta)-f'(\eta))^\top(\theta-\eta)$, and consequently
$(f'(\theta)-f'(\eta))^\top(\theta-\eta)\le L_f\norm{\theta-\eta}^2$.
\end{enumerate}
\end{lemma}

\begin{proof}
(b) By definition of $\mu$-strong convexity (Definition~5.2 of the book; the same computation with $\mu=0$,
i.e.\ with the definition of convexity, Definition~5.1, eq.~(5.6), gives (a)),
$f(\eta)\ge f(\theta)+f'(\theta)^\top(\eta-\theta)+\frac\mu2\norm{\eta-\theta}^2$ and
$f(\theta)\ge f(\eta)+f'(\eta)^\top(\theta-\eta)+\frac\mu2\norm{\eta-\theta}^2$. Summing the two
inequalities and cancelling $f(\theta)+f(\eta)$ gives
$0\ge -(f'(\theta)-f'(\eta))^\top(\theta-\eta)+\mu\norm{\theta-\eta}^2$, which is the claim.

(c) Fix $\theta$ and let $\varphi(\eta)=f(\eta)-f'(\theta)^\top\eta$. The function $\varphi$ is convex,
$L_f$-smooth, and $\varphi'(\theta)=0$, so $\theta$ minimizes $\varphi$. The descent lemma for
$L_f$-smooth functions (Definition~5.3 of the book: $\varphi(\eta')\le\varphi(\eta)+\varphi'(\eta)^\top(\eta'-\eta)+\frac{L_f}2\norm{\eta'-\eta}^2$)
applied with $\eta'=\eta-\frac1{L_f}\varphi'(\eta)$ gives
$\varphi(\theta)\le\varphi(\eta')\le\varphi(\eta)-\frac1{2L_f}\norm{\varphi'(\eta)}^2$, that is,
\[
f(\theta)\;\le\; f(\eta)+f'(\theta)^\top(\theta-\eta)-\frac1{2L_f}\norm{f'(\eta)-f'(\theta)}^2 .
\]
This is equation~(5.12) of the book. Writing it also with $\theta$ and $\eta$ exchanged and adding the two
inequalities yields $0\le (f'(\theta)-f'(\eta))^\top(\theta-\eta)-\frac1{L_f}\norm{f'(\theta)-f'(\eta)}^2$,
the co-coercivity inequality. The consequence follows from the Cauchy--Schwarz inequality
$(f'(\theta)-f'(\eta))^\top(\theta-\eta)\le\norm{f'(\theta)-f'(\eta)}\norm{\theta-\eta}$ together with the
Lipschitz property $\norm{f'(\theta)-f'(\eta)}\le L_f\norm{\theta-\eta}$.
\end{proof}

\begin{lemma}[$\muEff$ lies between $\mu$ and $L$]\label{lem:B2-mueff}
Under Assumption~\ref{ass:B2}, for all $\theta\neq\etastar$,
\begin{equation}\label{ineq:mueff-lower}
\mu\norm{\theta-\etastar}^2\;\le\;(\theta-\etastar)^\top\Fp(\theta)
\qquad\text{(strong convexity at the iterate; i.e.\ } \mu\le\muEff(\theta)),
\end{equation}
\begin{equation}\label{ineq:mueff-upper}
(\theta-\etastar)^\top\Fp(\theta)\;\le\;L\norm{\theta-\etastar}^2
\qquad\text{(smoothness at the iterate; i.e.\ } \muEff(\theta)\le L).
\end{equation}
\end{lemma}

\begin{proof}
Apply Lemma~\ref{lem:B2-twopoint}(b) and (c) to $f=F$ with $\eta=\etastar$ and use $\Fp(\etastar)=0$.
\end{proof}

The tightness ratio of \eqref{ineq:mueff-lower} at $\theta$ is exactly $\mu/\muEff(\theta)\in[\mu/L,1]$; this
is the quantity the experiments of the B3 note record under the name $r_{\mathrm{sc}}$ (Definition~\ref{def:B3-ratios}
of that note, which likewise calls $r_{\mathrm{coco}}=\LEff/R^2$ the tightness ratio of the co-coercivity step, see
Lemma~\ref{lem:B2-Leff}); we use the same two names below.

\begin{lemma}[integral form of the effective curvature]\label{lem:B2-integral}
Assume $F$ is twice continuously differentiable and let, for $\theta\neq\etastar$,
\[
\bar H(\theta)\;=\;\int_0^1 F''(\etastar+s\,v)\,ds \qquad (v=\theta-\etastar,\ u=v/\norm{v}).
\]
Then $\bar H(\theta)$ is symmetric, $\mu I\preccurlyeq\bar H(\theta)\preccurlyeq L I$, $\Fp(\theta)=\bar H(\theta)\,v$, and
\begin{equation}\label{eq:B2-rayleigh}
\muEff(\theta)\;=\;\frac{v^\top\bar H(\theta)\,v}{\norm{v}^2}\;=\;u^\top\bar H(\theta)\,u ,
\end{equation}
i.e.\ $\muEff(\theta)$ is the Rayleigh quotient, in the direction of the error, of the Hessian averaged along the
segment $[\etastar,\theta]$.
\end{lemma}

\begin{proof}
Let $g(s)=\Fp(\etastar+sv)$ for $s\in[0,1]$. Since $F$ is $C^2$, $g$ is $C^1$ with
$g'(s)=F''(\etastar+sv)\,v$, and the fundamental theorem of calculus gives
$\Fp(\theta)-\Fp(\etastar)=g(1)-g(0)=\int_0^1F''(\etastar+sv)v\,ds=\bar H(\theta)v$;
as $\Fp(\etastar)=0$, this is $\Fp(\theta)=\bar H(\theta)v$. The integral of a continuous function
with values in the set $\{A\ \text{symmetric}: \mu I\preccurlyeq A\preccurlyeq LI\}$ (a closed convex set,
since it is defined by linear inequalities $\mu\norm{w}^2\le w^\top Aw\le L\norm{w}^2$ for all $w$)
stays in that set, because Riemann sums are convex combinations; and for a twice differentiable
function, $\mu$-strong convexity and $L$-smoothness are equivalent to $\mu I\preccurlyeq F''(\theta)\preccurlyeq LI$
for all $\theta$ (Section~5.2.3 of \cite{bach2024learning}). Finally
$\muEff(\theta)=v^\top\Fp(\theta)/\norm{v}^2=v^\top\bar H(\theta)v/\norm{v}^2$.
\end{proof}

\begin{lemma}[$\LEff$ lies between $\muEff$ and $R^2$]\label{lem:B2-Leff}
Under Assumption~\ref{ass:B2}, for all $\theta\neq\etastar$, with $v=\theta-\etastar$,
\begin{equation}\label{ineq:Leff-jensen}
\norm{\Fp(\theta)}^2\;\le\;\frac1n\sum_{i=1}^n\norm{\fip(\theta)-\fip(\etastar)}^2
\qquad\text{(Jensen / gradient diversity)},
\end{equation}
\begin{equation}\label{ineq:Leff-cocoercive}
\frac1n\sum_{i=1}^n\norm{\fip(\theta)-\fip(\etastar)}^2\;\le\;R^2\,v^\top\Fp(\theta)
\qquad\text{(averaged co-coercivity, eq.~(5.31) of the book)},
\end{equation}
\begin{equation}\label{ineq:Leff-mu}
\mu\;\le\;\muEff(\theta)\;\le\;\frac{\norm{\Fp(\theta)}}{\norm{v}}\;\le\;\frac{\norm{\Fp(\theta)}^2}{v^\top\Fp(\theta)} .
\end{equation}
Consequently
\begin{equation}\label{eq:B2-Leff-sandwich}
\mu\;\le\;\muEff(\theta)\;\le\;\frac{\norm{\Fp(\theta)}^2}{v^\top\Fp(\theta)}\;\le\;\LEff(\theta)\;\le\;R^2 ,
\end{equation}
and, more precisely, with $h_i=\fip(\theta)-\fip(\etastar)$ (the vectors $a_i$ of Section~\ref{sec:B3-exact}) and
$\bar h=\frac1n\sum_ih_i=\Fp(\theta)$,
\begin{equation}\label{eq:B2-Leff-variance}
\LEff(\theta)\;=\;\frac{\norm{\Fp(\theta)}^2+\frac1n\sum_{i=1}^n\norm{h_i-\bar h}^2}{v^\top\Fp(\theta)} .
\end{equation}
If moreover each $f_i$ is $L_i$-smooth, then $\LEff(\theta)\le\max_iL_i$ (so one may take $R^2=\max_iL_i$).
\end{lemma}

\begin{proof}
\eqref{eq:B2-Leff-variance} is the variance decomposition
$\frac1n\sum_i\norm{h_i}^2=\norm{\bar h}^2+\frac1n\sum_i\norm{h_i-\bar h}^2$ (expand the squares and use
$\sum_i(h_i-\bar h)=0$), divided by $v^\top\Fp(\theta)>0$; dropping the nonnegative variance term gives
\eqref{ineq:Leff-jensen}. For \eqref{ineq:Leff-cocoercive}, apply Lemma~\ref{lem:B2-twopoint}(c) to each
$f_i$ (convex, $R^2$-smooth) with $\eta=\etastar$: $\norm{h_i}^2\le R^2h_i^\top v$; averaging over $i$ gives
$\frac1n\sum_i\norm{h_i}^2\le R^2\bar h^\top v=R^2v^\top\Fp(\theta)$. If $f_i$ is $L_i$-smooth, the same
argument gives $\frac1n\sum_i\norm{h_i}^2\le\frac1n\sum_iL_ih_i^\top v\le(\max_iL_i)\,\bar h^\top v$, using
$h_i^\top v\ge0$ (Lemma~\ref{lem:B2-twopoint}(a)). For \eqref{ineq:Leff-mu}: the first inequality is
Lemma~\ref{lem:B2-mueff}; the second is Cauchy--Schwarz, $v^\top\Fp(\theta)\le\norm{v}\norm{\Fp(\theta)}$,
divided by $\norm{v}^2$; the third is the same Cauchy--Schwarz inequality in the form
$\norm{\Fp(\theta)}/\norm{v}\le\norm{\Fp(\theta)}^2/(v^\top\Fp(\theta))$. Combining
\eqref{ineq:Leff-jensen}--\eqref{ineq:Leff-mu} with the definition of $\LEff$ gives \eqref{eq:B2-Leff-sandwich}.
\end{proof}

\begin{remark}
The chain \eqref{eq:B2-Leff-sandwich} shows that the ``effective condition number'' $\LEff(\theta)/\muEff(\theta)$
is always at least $1$, and that two effects are separated in $\LEff(\theta)$, whose ratio to $R^2$ is the tightness ratio
of the co-coercivity step:
\eqref{ineq:Leff-jensen} measures the \emph{diversity} of the individual gradient differences
(its tightness ratio $\norm{\Fp(\theta)}^2/\frac1n\sum_i\norm{h_i}^2$ equals $1$ iff all $h_i$ coincide),
while the remaining factor $\norm{\Fp(\theta)}^2/(v^\top\Fp(\theta))\le\LEff(\theta)\le R^2$ measures how much
of the bound $R^2$ is actually used by the full gradient in the direction of the error. Both ratios are cheap to
compute once $\etastar$ is known (one pass over the data).
\end{remark}

\begin{proposition}[directional limits at the optimum]\label{prop:B2-limits}
Assume each $f_i$ is twice continuously differentiable, and write $H_\ast=F''(\etastar)$ and
$M_\ast=\frac1n\sum_{i=1}^nf_i''(\etastar)^2$. For a unit vector $u$ define
\[
\Lambda(u)\;=\;\frac{u^\top M_\ast u}{u^\top H_\ast u}.
\]
Then, as $\theta\to\etastar$ (with $u=(\theta-\etastar)/\norm{\theta-\etastar}$ allowed to vary),
\[
\muEff(\theta)-u^\top H_\ast u\;\to\;0,\qquad \LEff(\theta)-\Lambda(u)\;\to\;0 ,
\]
uniformly in the direction. In particular, if $\theta\to\etastar$ along a fixed direction $u$, then
$\muEff(\theta)\to u^\top H_\ast u$ (the Rayleigh quotient of the Hessian at the optimum in the direction of
approach) and $\LEff(\theta)\to\Lambda(u)$. Moreover, for every unit $u$,
\begin{equation}\label{eq:B2-Lambda-bounds}
u^\top H_\ast u\;\le\;\frac{u^\top H_\ast^2u}{u^\top H_\ast u}\;\le\;\Lambda(u)\;\le\;R^2 ,
\end{equation}
consistently with \eqref{eq:B2-Leff-sandwich}.
\end{proposition}

\begin{proof}
Let $r=\norm{v}$ and $\bar H_i(\theta)=\int_0^1f_i''(\etastar+sv)\,ds$, so that, exactly as in
Lemma~\ref{lem:B2-integral}, $\fip(\theta)-\fip(\etastar)=\bar H_i(\theta)v$, $0\preccurlyeq\bar H_i(\theta)\preccurlyeq R^2I$,
and $\bar H(\theta)=\frac1n\sum_i\bar H_i(\theta)$. Set
$\varepsilon_i(r)=\sup_{\norm{w}\le r}\norm{f_i''(\etastar+w)-f_i''(\etastar)}_{\mathrm{op}}$ and
$\varepsilon(r)=\frac1n\sum_i\varepsilon_i(r)$; by continuity of $f_i''$ at $\etastar$, $\varepsilon(r)\to0$ as $r\to0$.
Since every point $\etastar+sv$ of the segment is within distance $r$ of $\etastar$,
$\norm{\bar H_i(\theta)-f_i''(\etastar)}_{\mathrm{op}}\le\varepsilon_i(r)$, hence
$\norm{\bar H(\theta)-H_\ast}_{\mathrm{op}}\le\varepsilon(r)$ and, by \eqref{eq:B2-rayleigh},
$\abs{\muEff(\theta)-u^\top H_\ast u}=\abs{u^\top(\bar H(\theta)-H_\ast)u}\le\varepsilon(r)$, which proves the first limit.

For the second, write $\LEff(\theta)=\mathcal N(\theta)/\mathcal D(\theta)$ with
$\mathcal N(\theta)=\frac1{r^2}\cdot\frac1n\sum_i\norm{\bar H_i(\theta)v}^2=\frac1n\sum_iu^\top\bar H_i(\theta)^2u$ and
$\mathcal D(\theta)=\frac1{r^2}v^\top\Fp(\theta)=u^\top\bar H(\theta)u=\muEff(\theta)$. For symmetric matrices
$0\preccurlyeq P,Q\preccurlyeq R^2I$ one has $P^2-Q^2=P(P-Q)+(P-Q)Q$, so
$\norm{P^2-Q^2}_{\mathrm{op}}\le2R^2\norm{P-Q}_{\mathrm{op}}$; therefore
$\abs{\mathcal N(\theta)-u^\top M_\ast u}\le2R^2\varepsilon(r)$, while $\abs{\mathcal D(\theta)-u^\top H_\ast u}\le\varepsilon(r)$.
Both denominators are at least $\mu$ (Lemma~\ref{lem:B2-mueff} and $H_\ast\succcurlyeq\mu I$) and both
numerators are at most $R^4$, so
\[
\abs{\LEff(\theta)-\Lambda(u)}
=\abs{\frac{\mathcal N}{\mathcal D}-\frac{\mathcal N_\ast}{\mathcal D_\ast}}
\le\frac{\abs{\mathcal N-\mathcal N_\ast}}{\mathcal D}+\mathcal N_\ast\frac{\abs{\mathcal D-\mathcal D_\ast}}{\mathcal D\,\mathcal D_\ast}
\le\frac{2R^2\varepsilon(r)}{\mu}+\frac{R^4\varepsilon(r)}{\mu^2}\;\xrightarrow[r\to0]{}\;0 ,
\]
uniformly in $u$ (here $\mathcal N_\ast=u^\top M_\ast u$, $\mathcal D_\ast=u^\top H_\ast u$).

For \eqref{eq:B2-Lambda-bounds}: with $A_i=f_i''(\etastar)$ and $\bar A=H_\ast$, the matrix
$\frac1n\sum_i(A_i-\bar A)^2=\frac1n\sum_iA_i^2-\bar A^2=M_\ast-H_\ast^2$ is positive semidefinite (each
$(A_i-\bar A)^2=(A_i-\bar A)^\top(A_i-\bar A)\succcurlyeq0$), which gives the middle inequality; the first one
is Cauchy--Schwarz, $(u^\top H_\ast u)^2\le\norm{H_\ast u}^2\norm{u}^2=u^\top H_\ast^2u$; the last one uses
$A_i^2\preccurlyeq R^2A_i$ for $0\preccurlyeq A_i\preccurlyeq R^2I$ (indeed $R^2A_i-A_i^2=A_i^{1/2}(R^2I-A_i)A_i^{1/2}\succcurlyeq0$),
so that $M_\ast\preccurlyeq R^2H_\ast$ and $\Lambda(u)\le R^2$.
\end{proof}

\begin{remark}
Proposition~\ref{prop:B2-limits} says that, close to the optimum, $\muEff$ and $\LEff$ are not numbers but
functions of the direction of the error. This is the reason why the direction in which an algorithm approaches
$\etastar$ (Section~\ref{sec:B2-direction}) matters for the observed contraction rate. The ``local
strong convexity constant'' $\lmin(H_\ast)$ of \cite{bach2014adaptivity} is the worst case of
$u^\top H_\ast u$ over directions, and the quantity $u^\top H_\ast^2u/(u^\top H_\ast u)$ in
\eqref{eq:B2-Lambda-bounds} is the Rayleigh-type quotient that governs gradient descent on a quadratic.
\end{remark}

\subsection{Logistic regression}\label{sec:B2-logistic}

\paragraph{Setting and constants.}
Let $\ell(a)=\log(1+e^{-a})$ and $\sigma(a)=1/(1+e^{-a})$. Then
\[
\ell'(a)=\sigma(a)-1=-\sigma(-a),\qquad
\ell''(a)=\sigma(a)\sigma(-a)=\sigma(a)\big(1-\sigma(a)\big)=\sigma'(a)\in(0,\tfrac14],
\]
\[
\ell'''(a)=\ell''(a)\big(1-2\sigma(a)\big),
\]
so that $\ell''$ is even, maximal at $0$ with $\ell''(0)=1/4$, and non-increasing on $[0,\infty)$ (because
$1-2\sigma(a)\le0$ for $a\ge0$). The data are $(x_i,y_i)\in\R^d\times\{-1,1\}$, $i=1,\dots,n$; we write
$R_x=\max_i\norm{x_i}$, $\hat\Sigma=\frac1n\sum_ix_ix_i^\top$ (empirical second-moment matrix of the features),
$a_i(\theta)=y_ix_i^\top\theta$, and
\[
f_i(\theta)=\ell(a_i(\theta))+\frac\lambda2\norm{\theta}^2,\qquad
\fip(\theta)=\ell'(a_i(\theta))\,y_ix_i+\lambda\theta,\qquad
f_i''(\theta)=\ell''(a_i(\theta))\,x_ix_i^\top+\lambda I ,
\]
using $y_i^2=1$. Hence
\begin{equation}\label{eq:B2-logistic-hessian}
F''(\theta)=\frac1n\sum_{i=1}^n\ell''\big(y_ix_i^\top\theta\big)\,x_ix_i^\top+\lambda I .
\end{equation}
The function $f_i$ is $L_i$-smooth with $L_i=\norm{x_i}^2/4+\lambda$ (the largest eigenvalue of
$f_i''(\theta)$, $\ell''(a_i(\theta))\norm{x_i}^2+\lambda$, is at most $L_i$, with equality when $a_i(\theta)=0$), so the
per-sample smoothness constant of the project is $R^2=\max_iL_i=R_x^2/4+\lambda$. The smoothness constant of
$F$ is $L=\lambda+\lmax(\hat\Sigma)/4$: indeed $F''(\theta)\preccurlyeq\frac14\hat\Sigma+\lambda I$ for all
$\theta$, with equality at $\theta=0$. Finally $F''(\theta)\succcurlyeq\lambda I$ for all $\theta$, so $F$ is
$\lambda$-strongly convex; Proposition~\ref{prop:B2-global} shows that $\mu=\lambda$ is the exact global constant.

\begin{lemma}[mean-value weights and explicit formulas]\label{lem:B2-logistic-formulas}
Let $\theta\neq\etastar$, $v=\theta-\etastar$, $u=v/\norm{v}$, $c_i=x_i^\top u$, $a_i=a_i(\theta)$,
$a_i^\ast=a_i(\etastar)$, and define the mean-value weights
\begin{equation}\label{eq:B2-weights}
w_i(\theta)\;=\;\int_0^1\ell''\big(a_i^\ast+s(a_i-a_i^\ast)\big)\,ds\;\in\;(0,\tfrac14] .
\end{equation}
By the fundamental theorem of calculus (see the proof), the weights have the closed form
\begin{equation}\label{eq:B2-weights-closed}
w_i(\theta)=\frac{\ell'(a_i)-\ell'(a_i^\ast)}{a_i-a_i^\ast}=\frac{\sigma(a_i)-\sigma(a_i^\ast)}{a_i-a_i^\ast}\quad\text{if }a_i\neq a_i^\ast,
\qquad w_i(\theta)=\ell''(a_i^\ast)\quad\text{if }a_i=a_i^\ast,
\end{equation}
so that the formulas below are exact and computable in $O(nd)$ operations without any quadrature (in floating point,
when $\abs{a_i-a_i^\ast}$ is below the square root of the machine precision, use $\ell''(\frac12(a_i+a_i^\ast))$, which
agrees with $w_i(\theta)$ up to $O(\abs{a_i-a_i^\ast}^2)$, to avoid cancellation; alternatively $\muEff$ and $\LEff$ can be
computed directly from Definition~\ref{def:B2-eff}). Then
\begin{equation}\label{eq:B2-logistic-graddiff}
\fip(\theta)-\fip(\etastar)=w_i(\theta)\,(x_i^\top v)\,x_i+\lambda v,\qquad
\bar H(\theta)=\frac1n\sum_{i=1}^nw_i(\theta)\,x_ix_i^\top+\lambda I ,
\end{equation}
and, with $A(\theta)=\frac1n\sum_iw_i(\theta)c_i^2$ and $B(\theta)=\frac1n\sum_iw_i(\theta)^2c_i^2\norm{x_i}^2$,
\begin{equation}\label{eq:B2-logistic-mueff}
\muEff(\theta)=\lambda+A(\theta)=\lambda+\frac1n\sum_{i=1}^nw_i(\theta)\,(x_i^\top u)^2 ,
\end{equation}
\begin{equation}\label{eq:B2-logistic-Leff}
\LEff(\theta)=\frac{\lambda^2+2\lambda A(\theta)+B(\theta)}{\lambda+A(\theta)}
=\muEff(\theta)+\frac{B(\theta)-A(\theta)^2}{\muEff(\theta)},\qquad B(\theta)\ge A(\theta)^2 .
\end{equation}
\end{lemma}

\begin{proof}
The map $s\mapsto\ell'(a_i^\ast+s(a_i-a_i^\ast))$ is $C^1$ with derivative
$\ell''(a_i^\ast+s(a_i-a_i^\ast))(a_i-a_i^\ast)$, so the fundamental theorem of calculus gives
$\ell'(a_i)-\ell'(a_i^\ast)=w_i(\theta)(a_i-a_i^\ast)=w_i(\theta)\,y_ix_i^\top v$, which is \eqref{eq:B2-weights-closed}
(when $a_i=a_i^\ast$ the integrand of \eqref{eq:B2-weights} is the constant $\ell''(a_i^\ast)$; $\ell'=\sigma-1$ gives the
second expression). Therefore
$\fip(\theta)-\fip(\etastar)=[\ell'(a_i)-\ell'(a_i^\ast)]y_ix_i+\lambda v=w_i(\theta)(x_i^\top v)x_i+\lambda v$
(using $y_i^2=1$). Averaging over $i$ gives $\Fp(\theta)=[\frac1n\sum_iw_i(\theta)x_ix_i^\top+\lambda I]v$; by
linearity of the integral this matrix is exactly $\int_0^1F''(\etastar+sv)ds$ computed from
\eqref{eq:B2-logistic-hessian}, i.e.\ the matrix $\bar H(\theta)$ of Lemma~\ref{lem:B2-integral}. Formula
\eqref{eq:B2-logistic-mueff} is the Rayleigh quotient \eqref{eq:B2-rayleigh}. For \eqref{eq:B2-logistic-Leff},
$\fip(\theta)-\fip(\etastar)=\norm{v}\,(w_ic_ix_i+\lambda u)$ and
$\norm{w_ic_ix_i+\lambda u}^2=w_i^2c_i^2\norm{x_i}^2+2\lambda w_ic_i^2+\lambda^2$ (because $x_i^\top u=c_i$ and
$\norm{u}=1$), so $\frac1n\sum_i\norm{\fip(\theta)-\fip(\etastar)}^2=\norm{v}^2(B+2\lambda A+\lambda^2)$, while
$v^\top\Fp(\theta)=\norm{v}^2(\lambda+A)$; dividing gives the first expression, and
$\lambda^2+2\lambda A+B=(\lambda+A)^2+(B-A^2)$ gives the second. Finally $\norm{x_i}^2\ge c_i^2$
(Cauchy--Schwarz) gives $B\ge\frac1n\sum_i(w_ic_i^2)^2\ge(\frac1n\sum_iw_ic_i^2)^2=A^2$ (Jensen).
\end{proof}

\begin{lemma}[two-sided bounds on the effective curvature]\label{lem:B2-logistic-bounds}
Let $\theta\neq\etastar$. For each $i$, let $m_i(\theta)$ be the smaller of $\ell''(a_i^\ast)$ and $\ell''(a_i)$, and let
$m(\theta)=\min_im_i(\theta)$. Then $m_i(\theta)\le w_i(\theta)\le\frac14$, and
\begin{equation}\label{ineq:logistic-lower}
\lambda+m(\theta)\,u^\top\hat\Sigma u\;\le\;\lambda+\frac1n\sum_{i=1}^nm_i(\theta)(x_i^\top u)^2\;\le\;\muEff(\theta),
\end{equation}
\begin{equation}\label{ineq:logistic-upper}
\muEff(\theta)\;\le\;\lambda+\frac14\,u^\top\hat\Sigma u\;\le\;\lambda+\frac14\lmax(\hat\Sigma)=L .
\end{equation}
In particular $\muEff(\theta)\ge\lambda+m(\theta)\lmin(\hat\Sigma)$, where
$m(\theta)=\min_i\min\{\sigma'(y_ix_i^\top\etastar),\sigma'(y_ix_i^\top\theta)\}$ is the minimum of $\ell''=\sigma'$
over the $n$ segments $[a_i^\ast,a_i]$ and is computable in $O(nd)$ operations.
\end{lemma}

\begin{proof}
Every point $a=a_i^\ast+s(a_i-a_i^\ast)$, $s\in[0,1]$, of the segment satisfies
$\abs a\le\max\{\abs{a_i^\ast},\abs{a_i}\}$; since $\ell''$ is even and non-increasing in $\abs a$, $\ell''$ on the
segment is at least $\min\{\ell''(a_i^\ast),\ell''(a_i)\}=m_i(\theta)$ and at most
$1/4$. Integrating over $s$ gives $m_i(\theta)\le w_i(\theta)\le1/4$, and inserting these bounds into
\eqref{eq:B2-logistic-mueff}, with $\frac1n\sum_i(x_i^\top u)^2=u^\top\hat\Sigma u\in[\lmin(\hat\Sigma),\lmax(\hat\Sigma)]$,
gives \eqref{ineq:logistic-lower}--\eqref{ineq:logistic-upper}.
\end{proof}

\paragraph{Global versus local strong convexity.}
The next proposition states precisely in which sense ``the global strong convexity constant is exactly
$\lambda$'' for logistic regression. Recall that, for a twice continuously differentiable $F$, $F$ is
$\mu$-strongly convex if and only if $\theta\mapsto F(\theta)-\frac\mu2\norm{\theta}^2$ is convex
(Exercise~5.5 of \cite{bach2024learning}), if and only if $F''(\theta)\succcurlyeq\mu I$ for all $\theta$, so that
the best (largest) strong convexity constant is
\[
\mu_{\mathrm{glob}}\;=\;\sup\{\mu\ge0:\ F\text{ is $\mu$-strongly convex}\}\;=\;\inf_{\theta\in\R^d}\lmin\big(F''(\theta)\big).
\]

\begin{proposition}[the global strong convexity constant is $\lambda$]\label{prop:B2-global}
For $\ell_2$-regularized logistic regression, let $I_0=\{i:x_i\neq0\}$ (if $I_0=\emptyset$ then $F''\equiv\lambda I$ and there is nothing to prove).
\begin{enumerate}[label=(\alph*)]
\item There exists a unit vector $\hat d\in\R^d$ with $x_i^\top\hat d\neq0$ for every $i\in I_0$; in fact Lebesgue-almost
every direction has this property.
\item For every such $\hat d$, with $\delta(\hat d)=\min_{i\in I_0}\abs{x_i^\top\hat d}>0$, and every $t\ge0$,
\begin{equation}\label{ineq:global-escape}
\lambda\;\le\;\lmin\big(F''(t\hat d)\big)\;\le\;\lambda+R_x^2\,e^{-t\,\delta(\hat d)} .
\end{equation}
\item Consequently $\mu_{\mathrm{glob}}=\inf_\theta\lmin(F''(\theta))=\lambda$: the global strong convexity
constant of $F$ is exactly $\lambda$, whatever the data.
\item The infimum is attained at some $\theta$ if and only if $\operatorname{span}\{x_1,\dots,x_n\}\neq\R^d$;
in that case $\lmin(F''(\theta))=\lambda$ for every $\theta$. If the features span $\R^d$, then
$\lmin(F''(\theta))>\lambda$ for every $\theta$ and the value $\lambda$ is only approached as
$\norm{\theta}\to\infty$; it is in fact approached along \emph{every} ray: for every unit vector $\hat d$ and every $t\ge0$,
\begin{equation}\label{ineq:global-ray}
\lmin\big(F''(t\hat d)\big)\;\le\;\lambda+\frac1n\sum_{i=1}^n\ell''\big(t\,y_ix_i^\top\hat d\big)(x_i^\top\hat d)^2\;\le\;\lambda+R_x^2\,e^{-t\,\delta'(\hat d)},
\end{equation}
where $\delta'(\hat d)=\min\{\abs{x_i^\top\hat d}:x_i^\top\hat d\neq0\}$, with the conventions $\delta'(\hat d)=+\infty$ and
$e^{-\infty}=0$ when $x_i^\top\hat d=0$ for all $i$; hence $\lim_{t\to\infty}\lmin(F''(t\hat d))=\lambda$ for every
direction $\hat d$. Directions as in (a) are those along which the whole matrix $F''(t\hat d)-\lambda I$ tends to $0$;
along a general ray,
\[
F''(t\hat d)-\lambda I\;\xrightarrow[t\to\infty]{}\;\frac1{4n}\sum_{i:\,x_i^\top\hat d=0}x_ix_i^\top ,
\]
a matrix which is singular because $\hat d$ lies in its kernel.
\end{enumerate}
\end{proposition}

\begin{proof}
(a) For $i\in I_0$ the set $\{\hat d:x_i^\top\hat d=0\}$ is a hyperplane, hence has Lebesgue measure zero; a finite
union of null sets is null, so its complement is nonempty (indeed of full measure); normalize any element.

(b) The lower bound is $F''\succcurlyeq\lambda I$. For the upper bound, $F''(t\hat d)-\lambda I=\frac1n\sum_{i\in I_0}
\ell''(t\,y_ix_i^\top\hat d)\,x_ix_i^\top$ (the terms with $x_i=0$ vanish), and this matrix is positive semidefinite with
\[
\norm{F''(t\hat d)-\lambda I}_{\mathrm{op}}\le\frac1n\sum_{i\in I_0}\ell''(t\,y_ix_i^\top\hat d)\norm{x_i}^2
\le R_x^2\max_{i\in I_0}\ell''(t\,y_ix_i^\top\hat d)\le R_x^2e^{-t\delta(\hat d)},
\]
where we used $\ell''(a)=\sigma(a)\sigma(-a)\le\sigma(-\abs a)=\frac1{1+e^{\abs a}}\le e^{-\abs a}$ (because
$\sigma\le1$ and $\ell''$ is even) and $\abs{t\,y_ix_i^\top\hat d}=t\abs{x_i^\top\hat d}\ge t\delta(\hat d)$. Hence
$\lmin(F''(t\hat d))\le\lambda+\norm{F''(t\hat d)-\lambda I}_{\mathrm{op}}\le\lambda+R_x^2e^{-t\delta(\hat d)}$.

(c) Letting $t\to\infty$ in \eqref{ineq:global-escape} gives $\inf_\theta\lmin(F''(\theta))\le\lambda$; the
reverse inequality is $F''\succcurlyeq\lambda I$.

(d) $\lmin(F''(\theta))=\lambda$ iff the positive semidefinite matrix $\frac1n\sum_i\ell''(a_i(\theta))x_ix_i^\top$ is
singular, iff there is $w\neq0$ with $\frac1n\sum_i\ell''(a_i(\theta))(x_i^\top w)^2=0$, iff $x_i^\top w=0$ for all
$i$ (since $\ell''>0$), iff the $x_i$ do not span $\R^d$. This condition does not depend on $\theta$.
For \eqref{ineq:global-ray}: $\lmin(F''(t\hat d))\le\hat d^\top F''(t\hat d)\,\hat d=\lambda+\frac1n\sum_i\ell''(t\,y_ix_i^\top\hat d)(x_i^\top\hat d)^2$
(Rayleigh quotient with $\norm{\hat d}=1$, using $y_i^2=1$); the terms with $x_i^\top\hat d=0$ vanish, and for the others
$\ell''(t\,y_ix_i^\top\hat d)(x_i^\top\hat d)^2\le e^{-t\abs{x_i^\top\hat d}}\norm{x_i}^2\le R_x^2e^{-t\delta'(\hat d)}$ by the bound
$\ell''(a)\le e^{-\abs a}$ used in (b). The limit of $F''(t\hat d)-\lambda I$ follows termwise from $\ell''(t\,y_ix_i^\top\hat d)\to0$
when $x_i^\top\hat d\neq0$ and $\ell''(0)=\frac14$ when $x_i^\top\hat d=0$.
\end{proof}

\begin{corollary}[strong convexity on a bounded region]\label{cor:B2-ball}
In $\ell_2$-regularized logistic regression, for every $D>0$ and every $\theta$ with $\norm\theta\le D$,
\begin{equation}\label{ineq:ball-lower}
\lmin\big(F''(\theta)\big)\;\ge\;\lambda+\ell''(R_xD)\,\lmin(\hat\Sigma)\;\ge\;\lambda+\tfrac14e^{-R_xD}\lmin(\hat\Sigma).
\end{equation}
Moreover $\norm{\etastar}\le\sqrt{2\log2/\lambda}$, so that the ball $\{\norm\theta\le2\sqrt{2\log2/\lambda}\}$
contains the ball of radius $\sqrt{2\log2/\lambda}$ around $\etastar$.
\end{corollary}

\begin{proof}
If $\norm\theta\le D$ then $\abs{a_i(\theta)}\le\norm{x_i}\norm\theta\le R_xD$, hence
$\ell''(a_i(\theta))\ge\ell''(R_xD)$ ($\ell''$ even, non-increasing in $\abs a$) and
$F''(\theta)\succcurlyeq\lambda I+\ell''(R_xD)\hat\Sigma$. The second inequality uses
$\ell''(a)=\sigma(\abs a)\sigma(-\abs a)\ge\frac12\cdot\frac1{1+e^{\abs a}}\ge\frac14e^{-\abs a}$.
For the norm of $\etastar$: $\frac\lambda2\norm{\etastar}^2\le F(\etastar)\le F(0)=\log2$.
\end{proof}

Thus the global constant $\lambda$ is only ``seen'' at infinity: on any bounded region the Hessian is
uniformly larger than $\lambda I$ by a data-dependent amount, and an algorithm whose iterates stay in a
bounded region never experiences the curvature $\lambda$ unless $\lmin(\hat\Sigma)=0$ (for instance when $d>n$;
Lemma~\ref{lem:B2-span} below identifies which trajectories then actually see it) or in directions where
$\hat\Sigma$ is nearly singular. Gradient descent with $\gamma\le1/L$ stays in the ball of radius
$\norm{\theta_0-\etastar}$ around $\etastar$: with $v=\theta_{t-1}-\etastar$,
$\norm{\theta_t-\etastar}^2=\norm v^2-2\gamma v^\top\Fp(\theta_{t-1})+\gamma^2\norm{\Fp(\theta_{t-1})}^2\le\norm v^2-\gamma(2-\gamma L)\,v^\top\Fp(\theta_{t-1})\le\norm v^2$
by co-coercivity of $F$ (Lemma~\ref{lem:B2-twopoint}(c) with $\eta=\etastar$) and $v^\top\Fp(\theta_{t-1})\ge0$.
For SAGA the statement is only probabilistic: under Proposition~5.9, $V_t\ge0$ and $\E[V_t\mid\mathcal F_{t-1}]\le V_{t-1}$,
so $(V_t)_{t\ge0}$ is a nonnegative supermartingale; it converges almost surely, hence $\sup_tV_t<\infty$ almost surely,
and Ville's maximal inequality gives $\Prob(\sup_{t\ge0}V_t\ge c)\le V_0/c$ for every $c>0$. Since
$\norm{\theta_t-\etastar}^2\le V_t$, the SAGA iterates are almost surely bounded; they are not, in general, confined to a
deterministic ball (for $n=1$ SAGA is gradient descent with $\gamma=1/(4R^2)\le1/L$, which stays in the ball of radius
$\norm{\theta_0-\etastar}$). For instance, take $n=2$, $d=1$, $x_1=1$, $x_2=0.1$, $y=(1,-1)$, $\lambda=10^{-3}$,
$\theta_0=0$, the step and initialization of Proposition~5.9, and cycles of $2\cdot10^4$ draws of the index $2$ followed
by one draw of the index $1$. While the index $2$ is repeated, the iterate converges to the point where
$f_2'(\theta)=-z_1^{(t)}$, the stored $z_1^{(t)}$ being stale, and $\abs{\theta_t-\etastar}$, initially $2.7$, equals
about $396$, $1294$, $2192$ and $3090$ after the cycles $1$, $3$, $5$ and $7$: it grows by about $(1-x_2)/\lambda=900$
every two cycles (a draw of the index $1$ at a point $\theta\ll0$ stores $z_1\approx-1+\lambda\theta$, and the next
phase brings the iterate to about $\abs\theta+(1-x_2)/\lambda$), and keeps growing if the phases are long enough. Since
every finite sequence of indices has positive probability, $\Prob(\sup_t\norm{\theta_t-\etastar}\ge c)>0$ for every
$c>0$. The next lemma, which is the scalar case of the generalized
self-concordance property of the logistic loss (\cite{bach2010self}, Lemma~1 and Proposition~1, eq.~(6);
Assumption (A2) in \cite{bach2014adaptivity}), turns this into a quantitative comparison between
$\muEff(\theta)$ and the curvature at the optimum.

\begin{lemma}[generalized self-concordance of the logistic loss]\label{lem:B2-selfconc}
For all $a,b\in\R$,
\[
e^{-\abs{b-a}}\,\ell''(a)\;\le\;\ell''(b)\;\le\;e^{\abs{b-a}}\,\ell''(a) .
\]
Consequently, for $\theta\neq\etastar$ and all $i$, with $p_i=\ell''(a_i^\ast)=\sigma'(y_ix_i^\top\etastar)$,
\begin{equation}\label{eq:B2-weights-selfconc}
e^{-\abs{x_i^\top v}}\,p_i\;\le\;w_i(\theta)\;\le\;e^{\abs{x_i^\top v}}\,p_i,
\qquad\text{so that}\qquad
e^{-R_x\norm v}p_i\le w_i(\theta)\le e^{R_x\norm v}p_i .
\end{equation}
\end{lemma}

\begin{proof}
Since $\ell''>0$ and $\ell'''=\ell''(1-2\sigma)$, the function $\log\ell''$ is differentiable with
$\abs{(\log\ell'')'(a)}=\abs{1-2\sigma(a)}\le1$; integrating between $a$ and $b$ gives
$\abs{\log\ell''(b)-\log\ell''(a)}\le\abs{b-a}$, which is the claim. For \eqref{eq:B2-weights-selfconc}, every point
of the segment $[a_i^\ast,a_i]$ is at distance at most $\abs{a_i-a_i^\ast}=\abs{y_ix_i^\top v}=\abs{x_i^\top v}\le R_x\norm v$
from $a_i^\ast$; apply the first claim with $a=a_i^\ast$ and integrate over $s\in[0,1]$ in \eqref{eq:B2-weights}.
\end{proof}

\begin{corollary}[effective curvature versus curvature at the optimum]\label{cor:B2-local}
For $\ell_2$-regularized logistic regression, let $p_i=\sigma'(y_ix_i^\top\etastar)$,
$H_\ast=F''(\etastar)=\frac1n\sum_ip_ix_ix_i^\top+\lambda I$ and $p_{\min}=\min_ip_i$.
For every $\theta\neq\etastar$, with $u=(\theta-\etastar)/\norm{\theta-\etastar}$,
\begin{equation}\label{ineq:selfconc-lower}
\lambda+e^{-R_x\norm{\theta-\etastar}}\,\big(u^\top H_\ast u-\lambda\big)\;\le\;\muEff(\theta),
\end{equation}
\begin{equation}\label{ineq:selfconc-upper}
\muEff(\theta)\;\le\;\lambda+e^{R_x\norm{\theta-\etastar}}\,\big(u^\top H_\ast u-\lambda\big),
\end{equation}
and
\begin{equation}\label{ineq:local-lower}
u^\top H_\ast u\;\ge\;\lambda+p_{\min}\,u^\top\hat\Sigma u\;\ge\;\lambda+p_{\min}\lmin(\hat\Sigma).
\end{equation}
In particular, $\liminf_{\theta\to\etastar}\muEff(\theta)\ge\lambda+p_{\min}\lmin(\hat\Sigma)$, and for
$\norm{\theta-\etastar}\le1/R_x$ one has $\muEff(\theta)\ge\lambda+e^{-1}p_{\min}\lmin(\hat\Sigma)$.
If moreover $\theta\in S=\operatorname{span}\{x_1,\dots,x_n\}$ (which holds along the whole trajectory of gradient
descent and of SAGA when $\theta_0\in S$, Lemma~\ref{lem:B2-span}), then $u\in S$, and $\lmin(\hat\Sigma)$ may be
replaced, in \eqref{ineq:local-lower} and in the two consequences, by $\lambda^+_{\min}(\hat\Sigma)$, the smallest
nonzero eigenvalue of $\hat\Sigma$ (which is positive even when $d>n$).
\end{corollary}

\begin{proof}
By \eqref{eq:B2-logistic-mueff} and \eqref{eq:B2-weights-selfconc},
$\muEff(\theta)-\lambda=\frac1n\sum_iw_i(\theta)c_i^2$ lies between $e^{\mp R_x\norm v}\frac1n\sum_ip_ic_i^2$,
and $\frac1n\sum_ip_ic_i^2=u^\top H_\ast u-\lambda$. Inequality \eqref{ineq:local-lower} is $p_i\ge p_{\min}$
and $\frac1n\sum_ic_i^2=u^\top\hat\Sigma u\ge\lmin(\hat\Sigma)$. For the last claim, $\etastar\in S$
(Lemma~\ref{lem:B2-span}(a)), so $v=\theta-\etastar\in S$ and $u\in S$; the kernel of the positive semidefinite matrix
$\hat\Sigma$ is $\{w:x_i^\top w=0\ \forall i\}=S^\perp$, so its range is $S$ and its restriction to $S$ has exactly the
nonzero eigenvalues of $\hat\Sigma$, whence $u^\top\hat\Sigma u\ge\lambda^+_{\min}(\hat\Sigma)$ for unit $u\in S$.
\end{proof}

\paragraph{How the data spectrum and $\lambda$ control the gap $\muEff/\mu$.}
Since $\mu=\lambda$ exactly (Proposition~\ref{prop:B2-global}), the ratio measuring the slack of the
strong-convexity step along a trajectory is $\muEff(\theta_t)/\lambda$. Corollary~\ref{cor:B2-local} shows
that, once $R_x\norm{\theta_t-\etastar}\lesssim1$ (for SAGA this takes a number of passes that grows with $\kappa$ in
the runs of Remark~\ref{rem:B2-P2-outcome}, where $n=1000$: the first recorded state with $R_x\norm{\theta_t-\etastar}\le1$
comes after $0.6$ to $7.1$ passes for $\kappa\le10^3$, after $10.6$ and $24.7$ passes for $\kappa=10^4$ with
$\alpha\in\{0,1\}$, and not within the $60$ passes of the run for $\kappa=10^4$, $\alpha=2$; key
\texttt{saga\_first\_pass\_Rx\_dist\_le\_1} of \texttt{results/B2\_summary.json}), this
ratio is within a factor $e^{\pm1}$ of
\begin{equation}\label{eq:B2-gap-direction}
\frac{u_t^\top H_\ast u_t}{\lambda}\;=\;1+\frac1\lambda\cdot\frac1n\sum_{i=1}^np_i\,(x_i^\top u_t)^2,
\qquad u_t=\frac{\theta_t-\etastar}{\norm{\theta_t-\etastar}},
\end{equation}
which, as $u_t$ ranges over the unit sphere, spans the interval
$[1+\lmin(H_\ast-\lambda I)/\lambda,\;1+\lmax(H_\ast-\lambda I)/\lambda]$, where the ``data Hessian''
$H_\ast-\lambda I=\frac1n\sum_ip_ix_ix_i^\top$ satisfies $p_{\min}\hat\Sigma\preccurlyeq H_\ast-\lambda I\preccurlyeq\frac14\hat\Sigma$.
Three regimes follow.
\begin{itemize}
\item \emph{Strong regularization,} $\lambda\gg\lmax(\hat\Sigma)/4$: the ratio is at most
$1+\lmax(\hat\Sigma)/(4\lambda)\approx1$ in every direction, the strong-convexity step is essentially tight,
but then $\kappa-1=R_x^2/(4\lambda)\ll R_x^2/\lmax(\hat\Sigma)\le n$ (as $\hat\Sigma\succcurlyeq\frac1nx_ix_i^\top$ for
every $i$), so that $\kappa\ll n$ and the problem lies outside the regime $\kappa\gg n$ of (A4).
\item \emph{Intermediate regularization,} $\lmin(H_\ast-\lambda I)\ll\lambda\ll\lmax(H_\ast-\lambda I)$ (the
typical situation with anisotropic features: a decaying spectrum of $\hat\Sigma$, or $d>n$, in which case
$\lmin(\hat\Sigma)=0$ and $\lmin(H_\ast)=\lambda$ exactly, the $\lambda$-eigenspace of $H_\ast$ being the orthogonal
complement $S^\perp$ of the feature span $S=\operatorname{span}\{x_1,\dots,x_n\}$; by Lemma~\ref{lem:B2-span} below
this value $\lambda$ is experienced only by trajectories having a component outside $S$, hence never when
$\theta_0\in S$, e.g.\ $\theta_0=0$, in which case the relevant bottom eigenvalue is
$\lmin(H_\ast|_S)=\lambda+\lambda^+_{\min}(H_\ast-\lambda I)$, with $\lambda^+_{\min}$ the smallest nonzero eigenvalue):
the gap is \emph{direction dependent}, close to $1$
along the bottom eigenvectors of $H_\ast$ and of order $\lmax(H_\ast)/\lambda\gg1$ along the top ones. Whether
the slack is observed is then decided by the direction of the error $u_t$, which is the object of
Section~\ref{sec:B2-direction}.
\item \emph{Weak regularization with well-conditioned features,} $\lambda\ll p_{\min}\lmin(\hat\Sigma)$: the
gap is at least $1+p_{\min}\lmin(\hat\Sigma)/\lambda\gg1$ in \emph{every} direction, and the
$\mu$-term of Proposition~5.9 is pessimistic, whatever the algorithm, by at least the factor
$1+e^{-1}p_{\min}\lmin(\hat\Sigma)/\lambda$ as soon as $R_x\norm{\theta_t-\etastar}\le1$ and by at least
$1+p_{\min}\lmin(\hat\Sigma)/\lambda$ in the limit $t\to\infty$ (Corollary~\ref{cor:B2-local}); far from $\etastar$
the data term carries the weight $e^{-R_x\norm{\theta_t-\etastar}}$ (or $m(\theta_t)$ in
Lemma~\ref{lem:B2-logistic-bounds}), which can be tiny, so the factor is not uniform along the whole trajectory.
Here $p_{\min}$ is controlled by the margins at the optimum: $p_i\ge\frac14e^{-\abs{a_i^\ast}}$ with
$\abs{a_i^\ast}\le R_x\norm{\etastar}\le R_x\sqrt{2\log2/\lambda}$ (Corollary~\ref{cor:B2-ball}), so $p_{\min}$ is
of order $1/4$ when $\lambda\gtrsim R_x^2$ but may be exponentially small when $\lambda$ is small and the data are
nearly separable; in the latter case $H_\ast-\lambda I$ is dominated by the points with $\abs{a_i^\ast}=O(1)$ (the
points near the decision boundary), and the relevant quantity is the spectrum of $\frac1n\sum_ip_ix_ix_i^\top$
rather than that of $\hat\Sigma$.
\end{itemize}
In all cases, the statement ``the gap is small when $\lambda\gg\lmin(H_\ast-\lambda I)$'' is correct only
for the worst direction (the bottom eigenvector of $H_\ast$); in a typical (e.g.\ isotropically random)
direction the ratio \eqref{eq:B2-gap-direction} is $1+\tr(H_\ast-\lambda I)/(d\lambda)$, which can be large even
when $\lambda\gg\lmin(H_\ast-\lambda I)$. The tightness ratios to record are therefore
$(\lambda+m(\theta_t)u_t^\top\hat\Sigma u_t)/\muEff(\theta_t)$ for \eqref{ineq:logistic-lower},
$\muEff(\theta_t)/(\lambda+\frac14u_t^\top\hat\Sigma u_t)$ for \eqref{ineq:logistic-upper}, and the two
self-concordance ratios of \eqref{ineq:selfconc-lower}--\eqref{ineq:selfconc-upper}, all of which only require
$\etastar$ (from a high-precision solver), one pass over the data and the eigen-decomposition of $H_\ast$.

\subsection{Direction dependence: gradient descent versus SAGA}\label{sec:B2-direction}

Proposition~\ref{prop:B2-limits} shows that near the optimum $\muEff(\theta_t)$ is the Rayleigh quotient of
$H_\ast=F''(\etastar)$ in the direction $u_t$ of the error. Which direction an algorithm selects is therefore
part of the answer to (B2). We treat the quadratic case exactly and use it as a model for the logistic case.

\paragraph{An exact invariant subspace: the feature span.}
Before the quadratic model, we record an exact constraint on the directions that gradient descent and SAGA can
explore for logistic regression (the argument only uses $\fip(\theta)-\lambda\theta\in\operatorname{span}\{x_i\}$, so it
applies verbatim to ridge regression and to every $\ell_2$-regularized generalized linear model
$f_i(\theta)=\varphi_i(x_i^\top\theta)+\frac\lambda2\norm\theta^2$). Below, $z_j^{(t)}$ is the SAGA memory,
$\omega_t$ the SAGA direction, and, for $i=1,\dots,n$,
\[
g_i^{(t+1)}\;=\;\Fp(\theta_t)-\fip(\theta_t)+z_i^{(t)}-\frac1n\sum_{j=1}^nz_j^{(t)}
\]
are the $n$ possible values of the estimator error $\xi_{t+1}=\Fp(\theta_t)-\omega_{t+1}$ at the state $(\theta_t,z^{(t)})$
(for the quadratic model of Lemma~\ref{lem:B2-saga-quadratic} below these are the vectors $g_i^{(t+1)}$ defined there), and
$C_{t+1}=\frac1n\sum_ig_i^{(t+1)}g_i^{(t+1)\top}$ is the conditional covariance of $\omega_{t+1}$ given the state.

\begin{lemma}[invariance of the feature span]\label{lem:B2-span}
Let $S=\operatorname{span}\{x_1,\dots,x_n\}\subseteq\R^d$, let $S^\perp$ be its orthogonal complement, and write
$\theta=\theta_S+\theta_\perp$ for the orthogonal decomposition of a vector $\theta$ (similarly $v_S,v_\perp$).
For $\ell_2$-regularized logistic regression:
\begin{enumerate}[label=(\alph*)]
\item $\etastar\in S$, and $\fip(\theta)-\lambda\theta\in S$ for every $\theta$ and every $i$; in particular
$\fip(\theta)_\perp=\lambda\theta_\perp$ for all $i$, and $\Fp(\theta)_\perp=\lambda\theta_\perp$.
\item (Invariance) If $\theta_0\in S$, then, for gradient descent and for SAGA with the initialization
$z_i^{(0)}=\fip(\theta_0)$ of Proposition~5.9 (or with $z_i^{(0)}=0$), all iterates $\theta_t$, all memories $z_i^{(t)}$,
all directions $\omega_t$ and all error vectors $g_i^{(t)}$ lie in $S$, for every $t$ and every realization of the
indices. The same holds for the quadratic model of Lemma~\ref{lem:B2-saga-quadratic} whenever $H_iS\subseteq S$ for
every $i$ and $v_0,e_1^{(0)},\dots,e_n^{(0)}\in S$.
\item (The complementary component is autonomous) For an arbitrary $\theta_0$, the $S^\perp$-component of the
error $v_t=\theta_t-\etastar$ evolves as if $F$ were the quadratic $\frac\lambda2\norm{\theta_\perp}^2$. For gradient
descent, $v_{t,\perp}=(1-\gamma\lambda)^tv_{0,\perp}$. For SAGA, let $s_j(t)$ be the time at which the current memory was
computed, $z_j^{(t)}=f'_j(\theta_{s_j(t)})$ (so $s_j(0)=0$, $s_{i(t)}(t)=t-1$ and $s_j(t)=s_j(t-1)$ for $j\neq i(t)$); then
\[
v_{t,\perp}=(1-\gamma\lambda)\,v_{t-1,\perp}-\gamma\lambda\Big(\frac1n\sum_{j=1}^nv_{s_j(t-1),\perp}-v_{s_{i(t)}(t-1),\perp}\Big),
\]
and the error vectors satisfy
\[
g^{(t+1)}_{i,\perp}=\lambda\Big(v_{s_i(t),\perp}-\frac1n\sum_{j=1}^nv_{s_j(t),\perp}\Big),
\]
so that $\E[v_{t,\perp}]=(1-\gamma\lambda)^tv_{0,\perp}$, and the noise injected by SAGA in a unit direction $q\in S^\perp$ is
\[
q^\top C_{t+1}\,q\;=\;\lambda^2\,\Var_j\big(q^\top v_{s_j(t)}\big)
\;:=\;\lambda^2\cdot\frac1n\sum_{i=1}^n\Big(q^\top v_{s_i(t)}-\frac1n\sum_{j=1}^nq^\top v_{s_j(t)}\Big)^2 ,
\]
the empirical variance, over the memory, of the $q$-component of the stored errors.
\end{enumerate}
\end{lemma}

\begin{proof}
(a) $\fip(\theta)=\ell'(a_i(\theta))\,y_ix_i+\lambda\theta$ and $\ell'(a_i(\theta))\,y_ix_i\in S$; averaging over $i$,
$\Fp(\theta)-\lambda\theta\in S$. At $\theta=\etastar$, $\Fp(\etastar)=0$ gives
$\lambda\etastar=-\frac1n\sum_i\ell'(a_i^\ast)\,y_ix_i\in S$, hence $\etastar\in S$ because $\lambda>0$. Projecting
$\fip(\theta)=[\fip(\theta)-\lambda\theta]+\lambda\theta$ onto $S^\perp$ gives $\fip(\theta)_\perp=\lambda\theta_\perp$.

(b) Induction on $t$. At $t=0$, $\theta_0\in S$ and $z_i^{(0)}=\fip(\theta_0)\in S$ by (a), since $\lambda\theta_0\in S$
(and trivially if $z_i^{(0)}=0$). If $\theta_{t-1}\in S$ and $z_j^{(t-1)}\in S$ for all $j$, then $\fip(\theta_{t-1})\in S$
for all $i$ by (a), hence $\omega_t=f'_{i(t)}(\theta_{t-1})+\frac1n\sum_jz_j^{(t-1)}-z_{i(t)}^{(t-1)}\in S$ (for gradient
descent, $\Fp(\theta_{t-1})\in S$), $\theta_t=\theta_{t-1}-\gamma\omega_t\in S$, the updated memory
$z_{i(t)}^{(t)}=f'_{i(t)}(\theta_{t-1})$ lies in $S$ and the other memories are unchanged. Finally each $g_i^{(t+1)}$ is a
linear combination of vectors of $S$. For the quadratic model, $\fip(\theta)-\fip(\etastar)=H_iv$ with $H_iS\subseteq S$,
and the same induction applies to $v_t$, $e_i^{(t)}$ and $g_i^{(t+1)}$ (Lemma~\ref{lem:B2-saga-quadratic}(a)).

(c) By (a) and $\etastar\in S$, $\fip(\theta)_\perp=\lambda\theta_\perp=\lambda(\theta-\etastar)_\perp$ for every $\theta$ and $i$.
For gradient descent, $v_{t,\perp}=v_{t-1,\perp}-\gamma\Fp(\theta_{t-1})_\perp=(1-\gamma\lambda)v_{t-1,\perp}$. For SAGA,
$z_j^{(t-1)}=f'_j(\theta_{s_j(t-1)})$ has $S^\perp$-component $\lambda v_{s_j(t-1),\perp}$, so the $S^\perp$-component of
$\omega_t$ is $\lambda v_{t-1,\perp}+\lambda\frac1n\sum_jv_{s_j(t-1),\perp}-\lambda v_{s_{i(t)}(t-1),\perp}$, and
$v_{t,\perp}=v_{t-1,\perp}-\gamma\,\omega_{t,\perp}$ is the displayed recursion. Since $i(t)$ is uniform and independent of
the past, the bracket has conditional mean zero, so $\E[v_{t,\perp}\mid\mathcal F_{t-1}]=(1-\gamma\lambda)v_{t-1,\perp}$ and
$\E[v_{t,\perp}]=(1-\gamma\lambda)^tv_{0,\perp}$. In $g_i^{(t+1)}$ the components $\Fp(\theta_t)_\perp$ and
$\fip(\theta_t)_\perp$ are both equal to $\lambda v_{t,\perp}$ and cancel, leaving
$z_{i,\perp}^{(t)}-\frac1n\sum_jz_{j,\perp}^{(t)}=\lambda(v_{s_i(t),\perp}-\frac1n\sum_jv_{s_j(t),\perp})$. For $q\in S^\perp$,
$q^\top g_i^{(t+1)}=q^\top g_{i,\perp}^{(t+1)}$, and $q^\top C_{t+1}q=\frac1n\sum_i(q^\top g_i^{(t+1)})^2$ gives the last formula.
\end{proof}

The lemma has two consequences for the measurements of this section. First, with the standard initialization
$\theta_0=0$ (or any $\theta_0\in S$) the error $v_t$ never leaves $S$, so only the restriction $H_\ast|_S$ of the Hessian
at the optimum to $S$ is explored; $S$ is $H_\ast$-invariant (as $H_\ast=\frac1n\sum_ip_ix_ix_i^\top+\lambda I$), and
when $d>n$ the $\lambda$-eigenspace $S^\perp$ of $H_\ast$ is exactly the part that is never explored. Second, for a
generic $\theta_0\notin S$ the $S^\perp$-component of the error is governed by the smallest possible curvature $\lambda$,
exactly as for the quadratic $\frac\lambda2\norm\theta^2$, for both algorithms.

\paragraph{Gradient descent on a quadratic.}
Let $F(\theta)=F(\etastar)+\frac12(\theta-\etastar)^\top H(\theta-\etastar)$ with $H$ symmetric positive definite,
eigenvalues $\lambda_1\ge\dots\ge\lambda_d=\lmin>0$ and an orthonormal basis of eigenvectors $q_1,\dots,q_d$, and let
$K_{\min}=\{k:\lambda_k=\lmin\}$. For a quadratic, $\bar H(\theta)=H$ for all $\theta$, so by
Lemma~\ref{lem:B2-integral} $\muEff(\theta)=v^\top Hv/\norm v^2$ exactly, and, along a trajectory $(\theta_t)$ with
$v_t=\theta_t-\etastar\neq0$, the spectral profile of the error
\begin{equation}\label{eq:B2-profile}
\pi_k(t)\;=\;\frac{(q_k^\top v_t)^2}{\norm{v_t}^2},\qquad k=1,\dots,d,\qquad\sum_k\pi_k(t)=1,
\end{equation}
determines it: $\muEff(\theta_t)=\sum_k\lambda_k\pi_k(t)$.

\begin{proposition}[alignment of gradient descent with the flattest direction]\label{prop:B2-gd}
Let $F(\theta)=F(\etastar)+\frac12(\theta-\etastar)^\top H(\theta-\etastar)$, where $H$ is symmetric positive
definite with eigenvalues $\lambda_1\ge\dots\ge\lambda_d=\lmin>0$ and orthonormal eigenvectors $q_1,\dots,q_d$; let
$E_{\min}$ be the eigenspace of $\lmin$, $P$ the orthogonal projector onto $E_{\min}$, and
$\lambda'=\min\{\lambda_k:\lambda_k>\lmin\}$ the second smallest distinct eigenvalue (if $H=\lmin I$, then $P=I$ and
the statement is trivial). Let $\theta_t=\theta_{t-1}-\gamma\Fp(\theta_{t-1})$ be gradient descent,
$v_t=\theta_t-\etastar$, and $\pi_k(t)$ the spectral profile \eqref{eq:B2-profile} of the error.
Let $0<\gamma\le1/\lambda_1$ and $v_0\neq0$; if $\gamma=1/\lambda_1$, assume moreover that $v_0$ does not lie in the
eigenspace of $\lambda_1$ (otherwise $\theta_t=\etastar$ for all $t\ge1$). Then $v_t\neq0$ for every $t$, and:
\begin{enumerate}[label=(\alph*)]
\item $v_t=(I-\gamma H)^tv_0$, i.e.\ $q_k^\top v_t=(1-\gamma\lambda_k)^t\,q_k^\top v_0$ for all $k$, and
$\muEff(\theta_t)=\sum_k\lambda_k\pi_k(t)$ with $\pi_k(t)\propto(1-\gamma\lambda_k)^{2t}(q_k^\top v_0)^2$.
\item (Monotonicity) $t\mapsto\muEff(\theta_t)$ is non-increasing.
\item (Alignment) If $Pv_0\neq0$ (true for Lebesgue-almost every $v_0$), then with
$\rho_{\mathrm{gap}}=(1-\gamma\lambda')/(1-\gamma\lmin)\in[0,1)$,
\[
\frac{\norm{v_t-Pv_t}}{\norm{Pv_t}}\;\le\;\rho_{\mathrm{gap}}^{\,t}\,\frac{\norm{v_0-Pv_0}}{\norm{Pv_0}},
\]
so that $v_t/\norm{v_t}$ converges to $E_{\min}$ at the geometric rate $\rho_{\mathrm{gap}}$, and
\begin{equation}\label{ineq:gd-align}
0\;\le\;\muEff(\theta_t)-\lmin\;\le\;(\lambda_1-\lmin)\,\rho_{\mathrm{gap}}^{\,2t}\,\frac{\norm{v_0-Pv_0}^2}{\norm{Pv_0}^2}
\;\xrightarrow[t\to\infty]{}\;0 .
\end{equation}
\item (Invariant subspaces) If $S$ is a linear subspace with $HS\subseteq S$ and $v_0\in S$, then $v_t\in S$ for all $t$,
and (a)--(c) hold with $H$ replaced by its restriction $H|_S$: the sums run over the eigenvectors $q_k\in S$ only,
$\lmin$ is replaced by $\lmin(H|_S)=\min\{\lambda_k:q_k\in S\}$, $E_{\min}$ and $P$ by the corresponding eigenspace
within $S$ and its orthogonal projector, $\lambda'$ by the second smallest distinct eigenvalue of $H|_S$, and the
hypothesis $Pv_0\neq0$ by its analogue for the restricted projector.
\end{enumerate}
\end{proposition}

\begin{proof}
(a) $\Fp(\theta)=H(\theta-\etastar)$, so $v_t=v_{t-1}-\gamma Hv_{t-1}=(I-\gamma H)v_{t-1}$, and
$q_k^\top(I-\gamma H)^tv_0=(1-\gamma\lambda_k)^tq_k^\top v_0$ since $Hq_k=\lambda_kq_k$. The formula for
$\muEff$ follows from $v_t^\top Hv_t=\sum_k\lambda_k(q_k^\top v_t)^2$. Since $1-\gamma\lambda_k>0$ unless
$\lambda_k=\lambda_1=1/\gamma$, the assumptions give $v_t\neq0$ for every $t$.

(b) Put $w_k(t)=(q_k^\top v_t)^2=(1-\gamma\lambda_k)^{2t}(q_k^\top v_0)^2\ge0$, $W(t)=\sum_kw_k(t)=\norm{v_t}^2>0$, and
$\varphi(\lambda)=(1-\gamma\lambda)^2$, so that $w_k(t+1)=\varphi(\lambda_k)w_k(t)$. On $[0,1/\gamma]\supset[\lmin,\lambda_1]$
the function $\varphi$ is non-increasing, hence $(\lambda_k-\lambda_j)(\varphi(\lambda_k)-\varphi(\lambda_j))\le0$ for all
$k,j$. Writing $\varphi_k=\varphi(\lambda_k)$,
\begin{align*}
W(t)W(t+1)\big[\muEff(\theta_{t+1})-\muEff(\theta_t)\big]
&=\sum_{k,j}w_kw_j\lambda_k\varphi_k-\sum_{k,j}w_kw_j\lambda_k\varphi_j\\
&=\frac12\sum_{k,j}w_kw_j(\lambda_k-\lambda_j)(\varphi_k-\varphi_j)\;\le\;0,
\end{align*}
where the middle equality is obtained by symmetrizing the sum in $(k,j)$ (Chebyshev's sum inequality).

(c) $\norm{Pv_t}^2=\sum_{k\in K_{\min}}(1-\gamma\lmin)^{2t}(q_k^\top v_0)^2=(1-\gamma\lmin)^{2t}\norm{Pv_0}^2$, where
$1-\gamma\lmin>0$ because $\gamma\le1/\lambda_1<1/\lmin$. For $k\notin K_{\min}$, $\lambda'\le\lambda_k\le\lambda_1\le1/\gamma$
gives $0\le1-\gamma\lambda_k\le1-\gamma\lambda'$, hence
$\norm{v_t-Pv_t}^2=\sum_{k\notin K_{\min}}(1-\gamma\lambda_k)^{2t}(q_k^\top v_0)^2\le(1-\gamma\lambda')^{2t}\norm{v_0-Pv_0}^2$.
Dividing gives the first claim. Then
$\muEff(\theta_t)-\lmin=\sum_k(\lambda_k-\lmin)\pi_k(t)\le(\lambda_1-\lmin)\sum_{k\notin K_{\min}}\pi_k(t)
=(\lambda_1-\lmin)\norm{v_t-Pv_t}^2/\norm{v_t}^2\le(\lambda_1-\lmin)\norm{v_t-Pv_t}^2/\norm{Pv_t}^2$,
using $\norm{v_t}\ge\norm{Pv_t}$; the lower bound $0$ is $\lambda_k\ge\lmin$. The set $\{v_0:Pv_0=0\}$ is a
proper linear subspace, hence Lebesgue-null.

(d) Since $H$ is symmetric and $HS\subseteq S$, also $HS^\perp\subseteq S^\perp$ (for $w\in S^\perp$ and $s\in S$,
$s^\top Hw=(Hs)^\top w=0$), so an orthonormal eigenbasis $(q_k)$ of $H$ can be chosen with every $q_k$ in $S$ or in
$S^\perp$. Then $q_k^\top v_0=0$ for $q_k\in S^\perp$, hence by (a) $q_k^\top v_t=0$ for all $t$ and $v_t\in S$; all the
sums in (a)--(c) run over the $k$ with $q_k\in S$ only, and the proofs apply verbatim to this restricted spectrum.
\end{proof}

\begin{remark}[scope of Proposition~\ref{prop:B2-gd}]\label{rem:B2-gd-scope}
(i) The alignment time is of order $1/\log(1/\rho_{\mathrm{gap}})\approx1/(\gamma(\lambda'-\lmin))$ iterations;
when the spectrum of $H$ is dense near $\lmin$ this can exceed the horizon of an experiment, in which case
$\muEff(\theta_t)$ decreases monotonically towards $\lmin$ without reaching it: monotonicity (b) is the
robust prediction, the limit (c) the asymptotic one. (ii) The case $Pv_0=0$ is not exceptional: by
Lemma~\ref{lem:B2-span}, for logistic regression with $d>n$ and $\theta_0\in S$ the bottom eigenspace $E_{\min}=S^\perp$
of $H_\ast$ is orthogonal to the whole trajectory, and in the quadratic model (d) then gives the alignment with the
flattest direction of $H|_S$, not with $E_{\min}$. Both statements hold in exact arithmetic. Rounding errors re-inject a
component along $E_{\min}$; that this component remains invisible, except possibly close to the precision floor, is
\emph{not proved} here: it is the heuristic of Remark~\ref{rem:B2-gd-rounding}. (iii) For
$\gamma\in(1/\lambda_1,2/\lambda_1)$ the contraction factors are $\abs{1-\gamma\lambda_k}$ and the largest one may
be attained at $\lambda_1$ (for $\gamma=2/(\lambda_1+\lmin)$ both extreme components contract by the same factor
$(\lambda_1-\lmin)/(\lambda_1+\lmin)$, the top one with alternating sign), so (b)--(c)
fail; the project uses $\gamma_{\mathrm{GD}}\le1/L=1/\lambda_1$, as in the book.
\end{remark}

\begin{remark}[rounding errors and the feature span: a heuristic]\label{rem:B2-gd-rounding}
This remark is a heuristic, not a proved statement. It concerns gradient descent in the situation of
Remark~\ref{rem:B2-gd-scope}(ii): logistic regression with $d>n$, $\theta_0\in S$ and $E_{\min}=S^\perp$. \emph{What is
proved} is the exact-arithmetic invariance of Lemma~\ref{lem:B2-span}(b), the alignment of Proposition~\ref{prop:B2-gd}(d)
for a quadratic, and the following bound. Write the computed iterates as
$\theta_t=\theta_{t-1}-\gamma\Fp(\theta_{t-1})+\delta_t$, where $\delta_t$ collects all the rounding errors of step $t$
(including those made in computing the gradient). Since $\Fp(\theta)_\perp=\lambda\theta_\perp$ for every $\theta$ and
$\etastar\in S$ (Lemma~\ref{lem:B2-span}(a)), the $S^\perp$-component of the error satisfies exactly
$v_{t,\perp}=(1-\gamma\lambda)v_{t-1,\perp}+\delta_{t,\perp}$; hence, for $\theta_0\in S$ and $0<\gamma\lambda\le1$ (true for
$\gamma\le1/L$),
\[
\norm{v_{t,\perp}}\;\le\;\sum_{s=1}^t(1-\gamma\lambda)^{t-s}\norm{\delta_s}\;\le\;\frac{\max_{s\le t}\norm{\delta_s}}{\gamma\lambda} .
\]
\emph{What is not proved} is the size of $\delta_t$. In the standard model of floating-point arithmetic it is of order
$\varepsilon_{\mathrm{mach}}\approx10^{-16}$ times the size of the quantities that are added ($\norm{\theta_{t-1}}$, and
$\gamma$ times the individual terms of the gradient), up to factors depending on $n$ and $d$; in particular it does not
decrease with the error $\norm{v_{t-1}}$. Under this model the component of the error along $E_{\min}$ stays below a level
of order $\varepsilon_{\mathrm{mach}}\max_s\norm{\theta_s}/(\gamma\lambda)$ (and is typically much smaller, since rounding
errors do not add up coherently), while the component in $S$ keeps decreasing geometrically. The former can therefore
become comparable with the error only when the error is within a factor of order $1/(\gamma\lambda)$ of the rounding
level; above that level one observes the behaviour inside $S$, that is, the alignment of
Proposition~\ref{prop:B2-gd}(d) with the bottom eigenspace of $H|_S$ in the quadratic model (for logistic regression, the
limit described, itself heuristically, in Remark~\ref{rem:B2-gd-heuristic}). Consistently, in the sanity check of
Remark~\ref{rem:B2-gd-heuristic} the $S^\perp$-component of the error stays at the rounding level $10^{-14}$. No proved
statement of this note depends on this remark; prediction (P1) below uses it only when $S\neq\R^d$ (e.g.\ $d>n$) and
$\theta_0\in S$, a situation that does not occur in the experiments of the project, where the features span $\R^d$
(Remark~\ref{rem:B2-P2-outcome}).
\end{remark}

\begin{remark}[gradient descent beyond quadratics: a heuristic]\label{rem:B2-gd-heuristic}
For a general $C^2$ function, Lemma~\ref{lem:B2-integral} gives
$v_t=(I-\gamma H_\ast)v_{t-1}-\gamma(\bar H(\theta_{t-1})-H_\ast)v_{t-1}$, and
$\norm{\bar H(\theta)-H_\ast}_{\mathrm{op}}\to0$ as $\theta\to\etastar$. Gradient descent is thus asymptotically
the linear recursion of Proposition~\ref{prop:B2-gd} with $H=H_\ast$, and we expect, for a generic $\theta_0$,
$v_t/\norm{v_t}\to E_{\min}(H_\ast)$ and $\muEff(\theta_t)\to\lmin(H_\ast)$ (for logistic regression,
$\lmin(H_\ast)=\lambda+\lmin(\frac1n\sum_ip_ix_ix_i^\top)$, which equals $\lambda$ when $d>n$). When
$\theta_0\in S=\operatorname{span}\{x_1,\dots,x_n\}$, e.g.\ $\theta_0=0$, Lemma~\ref{lem:B2-span} confines the whole
trajectory to $S$, which is $H_\ast$-invariant, and the expectation becomes, as in Proposition~\ref{prop:B2-gd}(d),
\[
\muEff(\theta_t)\;\to\;\lmin(H_\ast|_S)\;=\;\lambda+\lambda^+_{\min}\Big(\frac1n\sum_{i=1}^np_ix_ix_i^\top\Big),
\]
where $\lambda^+_{\min}$ denotes the smallest nonzero eigenvalue: the matrix $\frac1n\sum_ip_ix_ix_i^\top$ has range $S$
(all $p_i>0$) and its restriction to $S$ is positive definite. For $d>n$ this limit is strictly larger than
$\lmin(H_\ast)=\lambda$, and the value $\lambda$ is approached only from a $\theta_0\notin S$, through the component
$v_{t,\perp}=(1-\gamma\lambda)^tv_{0,\perp}$ of Lemma~\ref{lem:B2-span}(c), which then dominates the error for large $t$.
(Sanity check, reproduced by the script \texttt{experiments/check\_B2\_fp\_heuristic.py}, whose output is
\texttt{results/B2\_fp\_check\_summary.json}. Instance: $n=20$, $d=40$, $\lambda=10^{-4}$; with NumPy's generator
\texttt{default\_rng(0)} (seed $0$) one draws, in this order, a $20\times40$ standard Gaussian matrix whose rows,
normalized to unit norm, are the features $x_i$, a standard Gaussian $\beta\in\R^{40}$ and a standard Gaussian noise
$\epsilon\in\R^{20}$, and one sets $y_i=\operatorname{sign}(x_i^\top\beta+0.5\,\epsilon_i)$. Gradient descent with
$\gamma=1/L$, $L=\lambda+\lmax(\hat\Sigma)/4$, is run for $3000$ iterations. From $\theta_0=0$, $\muEff(\theta_t)$
decreases monotonically to $1.356\cdot10^{-4}$, against $\lambda+\lambda^+_{\min}=1.346\cdot10^{-4}$, while the
$S^\perp$-component of $\theta_t-\etastar$ stays at the rounding level (at most $1.6\cdot10^{-14}$); from a standard
Gaussian $\theta_0$, the next draw of the same generator, $\muEff(\theta_t)$ reaches $\lambda$ to within
$2.6\cdot10^{-5}$ in relative terms. The seeds $1$ and $2$ behave alike: monotone decrease, $\muEff(\theta_{3000})$
within $1\%$ of $\lambda+\lambda^+_{\min}$ and $S^\perp$-component below $2\cdot10^{-14}$ from $\theta_0=0$, relative gap
to $\lambda$ below $3\cdot10^{-5}$ from the Gaussian $\theta_0$.) We state this as
a heuristic: a rigorous proof requires the stable-manifold theory of the $C^1$ map $T(\theta)=\theta-\gamma\Fp(\theta)$
at its fixed point $\etastar$, whose Jacobian $I-\gamma H_\ast$ has eigenvalues in $[0,1)$, and is not needed
here. For logistic regression the linearization error is explicit: by \eqref{eq:B2-weights-selfconc},
$\abs{w_i(\theta)-p_i}\le(e^{R_x\norm v}-1)p_i$, hence for every unit vector $\psi$,
$\abs{\psi^\top(\bar H(\theta)-H_\ast)\psi}\le\frac1n\sum_i\abs{w_i(\theta)-p_i}(x_i^\top\psi)^2\le(e^{R_x\norm v}-1)\,\psi^\top(H_\ast-\lambda I)\psi$,
i.e.\ $\norm{\bar H(\theta)-H_\ast}_{\mathrm{op}}\le(e^{R_x\norm{\theta-\etastar}}-1)\norm{H_\ast-\lambda I}_{\mathrm{op}}$,
and likewise $\abs{\muEff(\theta)-u^\top H_\ast u}\le(e^{R_x\norm{\theta-\etastar}}-1)(u^\top H_\ast u-\lambda)$
(Corollary~\ref{cor:B2-local}). Once $R_x\norm{\theta_t-\etastar}\le0.1$, the quadratic model is accurate to
about $10\%$ in relative terms, and the measured $\muEff(\theta_t)$ can be compared with the ``quadratic proxy''
$\sum_k\lambda_k\pi_k(t)$ built from the eigenpairs $(\lambda_k,q_k)$ of $H_\ast$.
\end{remark}

\paragraph{SAGA on a quadratic: a linear random recursion.}
For quadratic components $f_i$ with Hessians $H_i$ (for instance ridge regression, $H_i=x_ix_i^\top+\lambda I$; logistic
regression is of this form to first order near $\etastar$ with $H_i=f_i''(\etastar)$),
$\fip(\theta)-\fip(\etastar)=H_iv$ exactly, and SAGA is a linear recursion with random coefficients.

\begin{lemma}[SAGA on a quadratic: exact moment recursions]\label{lem:B2-saga-quadratic}
Let $f_i(\theta)=\frac12\theta^\top H_i\theta-b_i^\top\theta+c_i$ with $H_i$ symmetric, $0\preccurlyeq H_i\preccurlyeq R^2I$,
and $H=\frac1n\sum_iH_i\succcurlyeq\mu I$. Consider SAGA with step $\gamma$ and memory $z_i^{(t)}$ as in the book, write
$v_t=\theta_t-\etastar$, and define the memory errors $e_i^{(t)}=z_i^{(t)}-\fip(\etastar)$ and
$\bar e^{(t)}=\frac1n\sum_ie_i^{(t)}$. Let $i(t)$ be uniform on $\{1,\dots,n\}$ and independent of the past
$\mathcal F_{t-1}=\sigma(i(1),\dots,i(t-1))$.
\begin{enumerate}[label=(\alph*)]
\item (State recursion) $v_t=(I-\gamma H)v_{t-1}+\gamma\xi_t$ with
\[
\xi_t=\big(e^{(t-1)}_{i(t)}-\bar e^{(t-1)}\big)-\big(H_{i(t)}-H\big)v_{t-1},\qquad \Ei[\xi_t]=0,
\]
and $e^{(t)}_{i(t)}=H_{i(t)}v_{t-1}$, $e^{(t)}_j=e^{(t-1)}_j$ for $j\neq i(t)$. Moreover $\xi_t=\Fp(\theta_{t-1})-\omega_t$
is minus the error of the SAGA gradient estimator $\omega_t$.
\item (Mean) $\E[v_t]=(I-\gamma H)^tv_0$: the mean error follows exactly the gradient descent recursion.
\item (Second moment) With $C_t=\Ei[\xi_t\xi_t^\top]=\frac1n\sum_{i=1}^ng_i^{(t)}g_i^{(t)\top}$, where
$g_i^{(t)}=(e_i^{(t-1)}-\bar e^{(t-1)})-(H_i-H)v_{t-1}$ (so that $C_t$ is the conditional covariance of
$\omega_t$ given $\mathcal F_{t-1}$),
\[
\E[v_tv_t^\top]=(I-\gamma H)\,\E[v_{t-1}v_{t-1}^\top]\,(I-\gamma H)+\gamma^2\E[C_t],
\]
and therefore, for every eigenpair $(\lambda_k,q_k)$ of $H$,
\begin{equation}\label{eq:B2-saga-spectral}
\E\big[(q_k^\top v_t)^2\big]=(1-\gamma\lambda_k)^{2t}(q_k^\top v_0)^2+\gamma^2\sum_{s=1}^t(1-\gamma\lambda_k)^{2(t-s)}\,q_k^\top\E[C_s]\,q_k .
\end{equation}
\end{enumerate}
\end{lemma}

\begin{proof}
(a) By definition $\omega_t=f'_{i(t)}(\theta_{t-1})+\frac1n\sum_jz_j^{(t-1)}-z_{i(t)}^{(t-1)}$. Write
$f'_{i(t)}(\theta_{t-1})=f'_{i(t)}(\etastar)+H_{i(t)}v_{t-1}$, $z_j^{(t-1)}=f'_j(\etastar)+e_j^{(t-1)}$ and
$\frac1n\sum_jf'_j(\etastar)=\Fp(\etastar)=0$. Then
$\omega_t=H_{i(t)}v_{t-1}+\bar e^{(t-1)}-e_{i(t)}^{(t-1)}=Hv_{t-1}-\xi_t$, and $v_t=v_{t-1}-\gamma\omega_t$ gives the
recursion. Since $i(t)$ is uniform and independent of $\mathcal F_{t-1}$, $\Ei[e_{i(t)}^{(t-1)}]=\bar e^{(t-1)}$ and
$\Ei[H_{i(t)}]=H$, so $\Ei[\xi_t]=0$. The memory update is $z_{i(t)}^{(t)}=f'_{i(t)}(\theta_{t-1})$, i.e.\
$e_{i(t)}^{(t)}=f'_{i(t)}(\theta_{t-1})-f'_{i(t)}(\etastar)=H_{i(t)}v_{t-1}$.

(b) Take expectations in (a): $\E[v_t]=(I-\gamma H)\E[v_{t-1}]+\gamma\E[\Ei[\xi_t]]=(I-\gamma H)\E[v_{t-1}]$.

(c) Expand $v_tv_t^\top=(I-\gamma H)v_{t-1}v_{t-1}^\top(I-\gamma H)+\gamma\big[(I-\gamma H)v_{t-1}\xi_t^\top+\xi_tv_{t-1}^\top(I-\gamma H)\big]+\gamma^2\xi_t\xi_t^\top$.
Since $v_{t-1}$ is $\mathcal F_{t-1}$-measurable and $\Ei[\xi_t]=0$, the conditional expectation of the bracket
vanishes; $\Ei[\xi_t\xi_t^\top]=\frac1n\sum_ig_i^{(t)}g_i^{(t)\top}=C_t$ because $\xi_t=g_{i(t)}^{(t)}$. Taking total
expectations gives the matrix recursion; multiplying it on both sides by $q_k$ and unrolling from $t$ down to $0$
gives \eqref{eq:B2-saga-spectral}.
\end{proof}

\paragraph{Reading of \eqref{eq:B2-saga-spectral} and a quasi-stationary heuristic.}
The first term of \eqref{eq:B2-saga-spectral} is the gradient-descent term of Proposition~\ref{prop:B2-gd},
which concentrates the error on $E_{\min}$. The second term is the gradient-estimator noise accumulated along
the trajectory: the noise injected at time $s$ in direction $q_k$ is damped by $(1-\gamma\lambda_k)^{2(t-s)}$, i.e.\
it is remembered for about $1/(2\gamma\lambda_k)$ iterations, longer in flat directions. The noise covariance
$C_s$ is itself a quadratic function of the state $(v_{s-1},e^{(s-1)})$, with
$\tr C_s=\frac1n\sum_i\norm{g_i^{(s)}}^2\le\frac2n\sum_i\norm{e_i^{(s-1)}}^2+\frac2n\sum_i\norm{H_iv_{s-1}}^2$ (the two
quantities bounded in Step~1 of the proof of Proposition~5.9), so it decays at the rate of the Lyapunov
function. To extract a prediction, assume (quasi-stationarity) that in the linear-convergence regime
$\E[C_s]\approx\rho^{\,s-t}\E[C_t]$ for $s\le t$, where $\rho$ is the contraction factor of the second moments.
For the proved factor $\rho=1-\min\{1/(3n),3\mu/(16R^2)\}$, in the regime $\kappa\gg n$ with
$\gamma=1/(4R^2)$ one has $1-\rho=\frac34\gamma\mu$ while $1-(1-\gamma\lambda_k)^2=\gamma\lambda_k(2-\gamma\lambda_k)\ge\gamma\lambda_k\ge\gamma\mu$,
so $\rho>(1-\gamma\lambda_k)^2$ for every $k$: the noise term decays more slowly than the gradient-descent term in
\emph{every} direction and dominates for large $t$ (the experiments contradict this conclusion for $\kappa=10^4=10n$ and
$\alpha\in\{0,1\}$, see Remark~\ref{rem:B2-P2-outcome}). For the observed factor $\rho_{\mathrm{obs}}$ this strict comparison
is not established; what is guaranteed is the weaker inequality $\rho_{\mathrm{obs}}\ge(1-\gamma\lambda_k)^2$ for every
$k$ with $q_k^\top v_0\neq0$, because by (b) and Jensen's inequality
$\E[(q_k^\top v_t)^2]\ge(q_k^\top\E v_t)^2=(1-\gamma\lambda_k)^{2t}(q_k^\top v_0)^2$, so that the second moment along $q_k$
cannot contract faster than $(1-\gamma\lambda_k)^2$ per step; in particular, if $Pv_0\neq0$,
$1-\rho_{\mathrm{obs}}\le2\gamma\lmin-\gamma^2\lmin^2$, which keeps the denominators
$2\lambda_k-(1-\rho_{\mathrm{obs}})/\gamma\ge2(\lambda_k-\lmin)+\gamma\lmin^2>0$ appearing below positive (with $\lmin$ read
as $\lmin(H|_S)$ in the situation of Proposition~\ref{prop:B2-gd}(d)). Summing the geometric series,
\begin{equation}\label{eq:B2-saga-heuristic}
\E\big[(q_k^\top v_t)^2\big]\;\approx\;\gamma^2\,\frac{c_k(t)}{1-(1-\gamma\lambda_k)^2/\rho}
\;\approx\;\frac{\gamma\,c_k(t)}{2\lambda_k-(1-\rho)/\gamma},
\qquad c_k(t)=q_k^\top\E[C_t]q_k ,
\end{equation}
using $(1-\gamma\lambda_k)^2\approx1-2\gamma\lambda_k$ and $\rho\approx1$. The predicted spectral profile of SAGA is
therefore
\begin{equation}\label{eq:B2-saga-profile}
\pi_k^{\mathrm{SAGA}}(t)\;\propto\;\frac{c_k(t)}{2\lambda_k-(1-\rho)/\gamma},\qquad
\muEff(\theta_t)\;\approx\;\frac{\sum_k\lambda_k\,c_k(t)/(2\lambda_k-(1-\rho)/\gamma)}{\sum_kc_k(t)/(2\lambda_k-(1-\rho)/\gamma)} ,
\end{equation}
a weighted mean of the eigenvalues in which direction $k$ receives the weight $c_k/(2\lambda_k-(1-\rho)/\gamma)$:
the noise re-injects \emph{all} directions it feeds, and the weight of a flat direction is larger only by the
factor $1/\lambda_k$, not by the exponentially growing factor $((1-\gamma\lmin)/(1-\gamma\lambda_k))^{2t}$ of
gradient descent. Two illustrative noise shapes: (i) isotropic noise, $c_k\equiv c$, gives
$\muEff\approx\big[\sum_k\lambda_k/(2\lambda_k-\frac34\mu)\big]/\big[\sum_k1/(2\lambda_k-\frac34\mu)\big]$; with one
eigenvalue equal to $\mu$ and $d-1$ eigenvalues $L\gg d\mu$ this is $\approx\mu\,(0.8+(d-1)/2)/0.8$, i.e.\ larger
than $\lmin=\mu$ by a factor of order $d$; (ii) ``data-shaped'' noise, $c_k\propto\lambda_k-\lambda$ (for
logistic or ridge regression started at $\theta_0\in S=\operatorname{span}\{x_1,\dots,x_n\}$, e.g.\ $\theta_0=0$, all
vectors $g_i^{(t)}$ lie in $S$ by Lemma~\ref{lem:B2-span}(b), so the noise has no energy in the pure-regularization
directions $S^\perp$; for a generic $\theta_0$ these directions carry, by Lemma~\ref{lem:B2-span}(c), the noise
$c_q=\lambda^2\Var_j(q^\top v_{s_j})$, $q\in S^\perp$, i.e.\ $\lambda^2$ times the spread of the $S^\perp$-errors stored in
the memory, which is negligible compared with the data-direction noise unless $\lambda$ is comparable to the data
curvature), gives weights $\propto(\lambda_k-\lambda)/(2\lambda_k-\frac34\mu)$, which
vanish on the directions with $\lambda_k\approx\lambda$ and are $\approx\frac12$ on all directions with
$\lambda_k\gg\lambda$: $\muEff(\theta_t)$ is then of the order of the \emph{average} curvature of the data
directions, far above $\lmin(H_\ast)=\lambda+\lmin(\frac1n\sum_ip_ix_ix_i^\top)$ (and above $\lmin(H_\ast|_S)$ when
$\theta_0\in S$).

\paragraph{Testable predictions and what to measure.}
Let $(\lambda_k,q_k)$ be the eigenpairs of $H_\ast=F''(\etastar)$ (with $\etastar$ from a high-precision solver),
and record along each run: the profile $\pi_k(t)$ of \eqref{eq:B2-profile}, its effective number of directions
$d_{\mathrm{eff}}(t)=1/\sum_k\pi_k(t)^2\in[1,d]$, the exact $\muEff(\theta_t)$ (formula
\eqref{eq:B2-logistic-mueff} with the closed-form weights \eqref{eq:B2-weights-closed}, or directly
Definition~\ref{def:B2-eff}) and its quadratic proxy $\sum_k\lambda_k\pi_k(t)$, and, for SAGA, the noise spectrum
$c_k(t)=q_k^\top C_tq_k$ where, for a general (non-quadratic) $F$,
\[
C_{t+1}=\frac1n\sum_{i=1}^ng_ig_i^\top,\qquad g_i=\Fp(\theta_t)-\fip(\theta_t)+z_i^{(t)}-\frac1n\sum_jz_j^{(t)} ,
\]
is the conditional covariance of the SAGA estimator $\omega_{t+1}$ over the $n$ possible draws of $i(t+1)$, given the
state $(\theta_t,z^{(t)})$ (Lemma~\ref{lem:B2-saga-quadratic}(c), Lemma~\ref{lem:B2-span}). Forming all the $c_k(t)$
costs $O(nd^2)$ per evaluation (the product $Q^\top G$ with $G=[g_1,\dots,g_n]\in\R^{d\times n}$ and
$Q=[q_1,\dots,q_d]$, or $n$ rank-one updates of a $d\times d$ matrix); it is $O(nd+d^2)$ only if a single direction
$q_k$ is monitored. Done once per pass, this is affordable for $d$ up to a few thousands. The initial point must be
part of the protocol, since it selects the subspace explored (Lemma~\ref{lem:B2-span}): run every experiment both
from $\theta_0=0\in S$ and from a generic $\theta_0$ (e.g.\ standard Gaussian), and compare $\muEff(\theta_t)$ with
$\lmin(H_\ast|_S)=\lambda+\lambda^+_{\min}(H_\ast-\lambda I)$ in the first case and with $\lmin(H_\ast)$ in the second
(the two coincide when the features span $\R^d$).
\begin{enumerate}[label=(P\arabic*)]
\item \emph{Gradient descent}: $\muEff(\theta_t)$ is non-increasing (exactly for a quadratic,
Proposition~\ref{prop:B2-gd}(b)) and converges to $\lmin(H_\ast)$ for a generic $\theta_0$, and to
$\lmin(H_\ast|_S)=\lambda+\lambda^+_{\min}(H_\ast-\lambda I)$ when $\theta_0\in S$, e.g.\ $\theta_0=0$
(Proposition~\ref{prop:B2-gd}(d) and Lemma~\ref{lem:B2-span}, completed by the heuristic
Remarks~\ref{rem:B2-gd-rounding} and~\ref{rem:B2-gd-heuristic}); under alignment,
$\pi(t)$ concentrates on the corresponding bottom eigenspace, so that $d_{\mathrm{eff}}(t)$ converges to a value in
$[1,m]$, $m$ being the multiplicity of that bottom eigenvalue: the limit is $1/\sum_{k\in K_{\min}}(\pi_k^\infty)^2$ with
$\pi_k^\infty=(q_k^\top Pv_0)^2/\norm{Pv_0}^2$, which equals $m$ only when $Pv_0$ is equidistributed over the $m$
eigenvectors (generically $m=1$ and $d_{\mathrm{eff}}(t)\to1$); the convergence speed is governed by the spectral gap
$\lambda'-\lmin$ of the explored spectrum.
\item \emph{SAGA}: after the initial transient, $\pi_k(t)$ remains spread over the directions fed by the
noise, with $\pi_k(t)\propto c_k(t)/(2\lambda_k-(1-\rho)/\gamma)$ up to fluctuations \eqref{eq:B2-saga-profile};
$d_{\mathrm{eff}}(t)$ stays of the order of the number of data directions; $\muEff(\theta_t)$ fluctuates
around a level strictly above $\lmin(H_\ast)$ (resp.\ $\lmin(H_\ast|_S)$ when $\theta_0\in S$) and does not decrease
monotonically; for the same problem and
large $t$, $\muEff^{\mathrm{SAGA}}(t)>\muEff^{\mathrm{GD}}(t)$. The experiments confirm (P2) only for $\kappa$ at most of the
order of $n$, and not in every such run (Remark~\ref{rem:B2-P2-outcome}).
\item \emph{Consequence for (B2)}: the slack of the strong-convexity step \eqref{ineq:mueff-lower}, measured by
$\muEff(\theta_t)/\mu$, is larger for SAGA than the asymptotic slack $\lmin(H_\ast)/\mu$ (resp.\
$\lmin(H_\ast|_S)/\mu$ for $\theta_0\in S$) of gradient descent; in the ``intermediate regularization'' regime of
Section~\ref{sec:B2-logistic} this is precisely the situation where the data anisotropy makes the bottom eigenvalue of
the explored spectrum close to $\lambda$ while the trajectory-averaged curvature is much larger. A violation of (P2) (e.g.\ $\pi(t)$ concentrating on $E_{\min}$ for SAGA as well) would mean that
the noise term in \eqref{eq:B2-saga-spectral} decays faster than assumed, which can be checked directly on
$c_k(t)$.
\end{enumerate}

\begin{remark}[outcome of prediction (P2)]\label{rem:B2-P2-outcome}
\emph{Measurements} (Figure~B2 of the report shows the case ${\alpha=1}$). $\ell_2$-regularized logistic regression with
$n=1000$, $d=50$, Gaussian features with covariance spectrum $j^{-\alpha}$ (randomly rotated, rescaled so that
$R_x=\max_i\norm{x_i}=1$), $\alpha\in\{0,1,2\}$, $\kappa=R^2/\mu\in\{10^2,10^3,10^4\}$, a single seed (seed $0$ for the data
and for the sampling); $\theta_0=0$; SAGA with $\gamma=1/(4R^2)$ and $z_i^{(0)}=\fip(\theta_0)$ for $60$ passes, recorded every
$n/10$ iterations; gradient descent with $\gamma=1/L$ for $400$ iterations, recorded at every iteration. States with
$\norm{\theta_t-\etastar}^2<10^{-18}$ are discarded and ``end'' denotes the last recorded state above this floor (SAGA
reaches it after $34$ to $49$ passes for $\kappa=10^2$, and for $\kappa=10^3$ with $\alpha\in\{0,1\}$; the other SAGA runs
end after $60$ passes). Since $\lmin(H_\ast)>\lambda$ in every configuration (table below), the features span $\R^d$
(Proposition~\ref{prop:B2-global}(d)), so that $S=\R^d$ and $\lmin(H_\ast|_S)=\lmin(H_\ast)$. We write $q_{\min}$ for a unit
eigenvector of $H_\ast$ for its smallest eigenvalue (which is simple in all nine configurations),
$\pi_{\min}(t)=(q_{\min}^\top(\theta_t-\etastar))^2/\norm{\theta_t-\etastar}^2$ for its share in the profile
\eqref{eq:B2-profile}, and $\lambda'$ for the second smallest eigenvalue of $H_\ast$; ``g.\ mean'' is the geometric mean over
the recorded states. Gradient descent reaches the floor after $23$ to $171$ iterations for $\kappa=10^2$ and for
$(\alpha,\kappa)=(0,10^3)$, and otherwise ends at iteration $400$, with $\norm{\theta_t-\etastar}^2=1.5\cdot10^{-18}$,
$4.7\cdot10^{-9}$, $1.0$, $5.7\cdot10^{-4}$ and $202$ for $(\alpha,\kappa)=(0,10^4)$, $(1,10^3)$, $(1,10^4)$, $(2,10^3)$,
$(2,10^4)$ (from $\norm{\theta_0-\etastar}^2=486$, $1954$, $389$ and $2565$ for the last four). Its end values coincide
with the limit $\lmin(H_\ast)$ of (P1) only for $\alpha=0$, $\kappa\ge10^3$: for $\kappa=10^2$ it reaches the
floor before its error has aligned with $q_{\min}$ ($\pi_{\min}=0.53$, $0.074$ and $0.001$ at the end), for $\kappa=10^3$,
$\alpha\in\{1,2\}$ the error is not yet aligned ($\muEff/\mu$ ends $9\%$ and $11\%$ above $\lmin(H_\ast)/\mu$), and for
$\kappa=10^4$ it ends at the edge of the local regime ($\alpha=1$) or far from it ($\alpha=2$). The table and the other
numbers of this remark are read from, or computed from the values stored in, \texttt{results/B2\_summary.json}, written
by \texttt{experiments/exp\_B2\_mu\_eff.py} from the cached trajectories.
\begin{center}
\footnotesize
\setlength{\tabcolsep}{4pt}
\begin{tabular}{@{}cccccccccccc@{}}
\toprule
 & & & & \multicolumn{3}{c}{SAGA: $\muEff/\mu$} & GD: $\muEff/\mu$ & \multicolumn{2}{c}{$\pi_{\min}$ at the end} & \multicolumn{2}{c}{$d_{\mathrm{eff}}$ at the end}\\
\cmidrule(lr){5-7}\cmidrule(lr){8-8}\cmidrule(lr){9-10}\cmidrule(l){11-12}
$\alpha$ & $\kappa$ & $\lmin(H_\ast)/\mu$ & $9n\lmin(H_\ast)/(16R^2)$ & start & g.\ mean & end & end & SAGA & GD & SAGA & GD\\
\midrule
0 & $10^2$ & 1.609 & 9.05 & 2.086 & 1.996 & 2.088 & 1.635 & 0.0008 & 0.530 & 23.7 & 2.51\\
0 & $10^3$ & 3.275 & 1.84 & 7.917 & 7.137 & 7.341 & 3.275 & 0.015 & 1.000 & 23.7 & 1.000\\
0 & $10^4$ & 9.863 & 0.555 & 50.05 & 10.59 & 9.942 & 9.863 & 0.998 & 1.000 & 1.005 & 1.000\\
1 & $10^2$ & 1.100 & 6.19 & 1.793 & 1.901 & 2.147 & 1.155 & 0.0008 & 0.074 & 15.2 & 10.1\\
1 & $10^3$ & 1.737 & 0.977 & 4.632 & 2.198 & 1.904 & 1.900 & 0.308 & 0.148 & 3.79 & 3.89\\
1 & $10^4$ & 4.970 & 0.280 & 21.02 & 5.386 & 4.974 & 5.058 & 0.9996 & 0.996 & 1.001 & 1.008\\
2 & $10^2$ & 1.003 & 5.64 & 2.497 & 2.167 & 4.980 & 1.029 & $8\cdot10^{-6}$ & 0.0010 & 2.71 & 11.6\\
2 & $10^3$ & 1.020 & 0.573 & 5.340 & 1.163 & 1.062 & 1.129 & 0.0025 & 0.0010 & 12.1 & 14.6\\
2 & $10^4$ & 1.142 & 0.0643 & 13.79 & 1.809 & 1.468 & 1.787 & 0.0017 & 0.0007 & 11.9 & 15.3\\
\bottomrule
\end{tabular}
\end{center}
\emph{Outcome.} (P2) holds only for $\kappa$ at most of the order of $n$, and not in every such run: it is confirmed in
the four runs with $9n\lmin(H_\ast)/(16R^2)>1$ ($\kappa=10^2$, all $\alpha$; $\kappa=10^3$, $\alpha=0$), refuted for
$\kappa=10^4$ with $\alpha\in\{0,1\}$ (one seed each; the initial error already carries $89$--$98\%$ of its energy on
$q_{\min}$), and the three remaining runs ($\kappa=10^3$ with $\alpha\in\{1,2\}$; $\kappa=10^4$ with $\alpha=2$) are
intermediate, see (c).
\begin{enumerate}[label=(\alph*),leftmargin=*,topsep=2pt,itemsep=1pt]
\item For $\kappa=10^2$ (all $\alpha$) and for $\kappa=10^3$ with $\alpha=0$, the error of SAGA does not concentrate on
$q_{\min}$ ($\pi_{\min}\le0.05$ at every state of the last five recorded passes), and over these passes the geometric
mean of $\muEff/\mu$ is between $1.25$ and $2.6$ times $\lmin(H_\ast)/\mu$ (single states fluctuate between $1.2$ and
$6.0$ times; the ``end'' columns of the table are single states), whereas gradient descent ends closer to
$\lmin(H_\ast)/\mu$. For $\alpha\in\{0,1\}$ the error stays spread over many directions ($d_{\mathrm{eff}}$ averages $14$
to $18$ over the last five passes). The run $\kappa=10^3$, $\alpha=0$ is the clearest: the initial error
$-\etastar$ is concentrated on $q_{\min}$ ($\pi_{\min}(0)=0.96$), gradient descent keeps it there and SAGA spreads it out,
although the spectral gap is large ($\lambda'/\lmin(H_\ast)=1.46$). For $\alpha=2$, $\kappa=10^2$, $d_{\mathrm{eff}}$ is
smaller (mean $5.4$ over the last five passes), but the five lowest eigenvectors of $H_\ast$ carry only $0.005$ of the
error and the geometric mean of $\muEff/\mu$ is $2.6$: the error sits on intermediate- and high-curvature directions,
consistently with the ``data-shaped'' noise (ii) above, whose weights vanish on the directions with
$\lambda_k\approx\lambda$.
\item For $\kappa=10^4$ and $\alpha\in\{0,1\}$, the error of SAGA concentrates on $q_{\min}$ as that of gradient descent
does ($\pi_{\min}=0.998$ and $0.9996$, $d_{\mathrm{eff}}=1.005$ and $1.001$), and $\muEff/\mu$ ends within $1\%$ of
$\lmin(H_\ast)/\mu$. The initial error is already largely aligned ($\pi_{\min}(0)=0.98$ and $0.89$); SAGA keeps and sharpens
this alignment instead of spreading the error. The large initial values $\muEff(\theta_0)=50.05\,\mu$ and
$21.02\,\mu$ are far from the Rayleigh quotients of $H_\ast$ in the initial direction ($10.38\,\mu$ and $5.885\,\mu$); they
are Rayleigh quotients of the averaged Hessian $\bar H(\theta_0)$ of Lemma~\ref{lem:B2-integral}: the point
$\theta_0=0$ is far from $\etastar$ ($R_x\norm{\theta_0-\etastar}=32$ and $44$), where the comparison of
Corollary~\ref{cor:B2-local} is uninformative. (P3) fails accordingly: at the end of the run the slack $\muEff/\mu$ of SAGA
exceeds $\lmin(H_\ast)/\mu$ by less than $1\%$.
\item Intermediate cases. For $\kappa=10^3$, $\alpha=1$, the alignment with $q_{\min}$ is partial ($\pi_{\min}=0.31$,
$d_{\mathrm{eff}}=3.79$), but the five lowest eigenvectors of $H_\ast$ carry $0.88$ of the error at the end, and
$\lambda'/\lmin(H_\ast)=1.005$, so that alignment with $q_{\min}$ itself is slow (Remark~\ref{rem:B2-gd-scope}(i)). For
$\alpha=2$ and $\kappa\ge10^3$ the bottom of the spectrum is nearly degenerate ($29$ and $11$ eigenvalues of $H_\ast$ lie
within $10\%$ of $\lmin(H_\ast)$ for $\kappa=10^3$ and $10^4$) and neither algorithm aligns with $q_{\min}$
(for gradient descent, alignment would take of the order of $L/(\lambda'-\lmin(H_\ast))\approx9\cdot10^4$ iterations, far
more than the $400$ of the protocol, and at $\kappa=10^4$ it is moreover still far from $\etastar$); at $\kappa=10^3$
SAGA nevertheless ends at $\muEff/\mu=1.062$, $4.2\%$ above $\lmin(H_\ast)/\mu$, while at $\kappa=10^4$ the run is not
finished after $60$ passes: $\norm{\theta_t-\etastar}^2=8.6$ at the end ($R_x\norm{\theta_t-\etastar}\approx2.9$, outside
the region where the bounds of Corollary~\ref{cor:B2-local} are informative, although the quadratic proxy
$\sum_k\lambda_k\pi_k(t)$ already matches $\muEff(\theta_t)$ to $0.2\%$) and $\muEff/\mu$ is still decreasing.
\end{enumerate}

\emph{Explanation (a heuristic).} Write $\lmin=\lmin(H_\ast)$. Formula \eqref{eq:B2-saga-profile} involves the actual decay
factor $\rho$ of the noise covariance, and the spread predicted in (P2) came from inserting the certified factor of
Proposition~5.9, which exceeds $(1-\gamma\lambda_k)^2$ for every $k$ (its distance to $1$ is at most
$\frac34\gamma\mu<\gamma\lambda_k(2-\gamma\lambda_k)$). If instead the noise decays as fast as the slowest mode of the
mean dynamics, i.e.\ if $1-\rho$ is close to its largest admissible value $2\gamma\lmin-\gamma^2\lmin^2$ (discussion before
\eqref{eq:B2-saga-heuristic}), the denominator $2\lmin-(1-\rho)/\gamma$ of the weight of $q_{\min}$ becomes small compared
with those of the other directions, which are at least $2(\lambda_k-\lmin)$, and the same formula predicts concentration on
$q_{\min}$ (unless $q_{\min}^\top\E[C_t]q_{\min}$ is comparatively negligible). Which case occurs is decided by the memory.
For SAGA with $z_i^{(0)}=\fip(\theta_0)$ one has $z_j^{(t-1)}=f'_j(\theta_{s_j(t-1)})$, with $s_j$ as in
Lemma~\ref{lem:B2-span}(c). Writing $\zeta_j=f'_j(\theta_{s_j(t-1)})-f'_j(\theta_{t-1})$ and inserting
$z_j^{(t-1)}=f'_j(\theta_{t-1})+\zeta_j$ in the definition of $\omega_t$ gives, exactly and for any $F$,
\begin{align*}
\xi_t&=\Fp(\theta_{t-1})-\omega_t=\zeta_{i(t)}-\frac1n\sum_{j=1}^n\zeta_j,\\
\Ei\norm{\xi_t}^2&\le\frac1n\sum_{j=1}^n\norm{\zeta_j}^2\le\frac{R^4}n\sum_{j=1}^n\norm{\theta_{s_j(t-1)}-\theta_{t-1}}^2
\end{align*}
(a variance is at most the second moment, and the $\fip$ are $R^2$-Lipschitz): the noise is produced by the displacement of
the iterate since each memory slot was last refreshed, about $n$ iterations earlier on average. Along $q_{\min}$ the mean
error decays by the factor $(1-\gamma\lmin)^n\approx e^{-n\gamma\lmin}$ per pass (Lemma~\ref{lem:B2-saga-quadratic}(b)), and
faster along the other eigenvectors. If $n\gamma\lmin$ is small, then, once the faster components have decayed, the iterate
moves little during one refresh cycle (in relative terms, by about $n\gamma\lmin$ along $q_{\min}$), the noise is a small
perturbation of the gradient-descent dynamics and decays with the slow component, and the error aligns with $q_{\min}$. If
$n\gamma\lmin\gtrsim1$, every mode decays by a constant factor within one refresh cycle, the stored gradients were computed
much farther from $\etastar$ than $\theta_{t-1}$, the noise decays only as fast as the memory is refreshed, more slowly than
every mode, and the error stays spread. For $\gamma=1/(4R^2)$ the condition $\frac94n\gamma\lmin<1$ is exactly
$\frac{3\lmin}{16R^2}<\frac1{3n}$: up to a constant factor, the dividing line is the condition that the parameter term be
the active one in the minimum of Proposition~5.9 with $\mu$ replaced by the local constant $\lmin(H_\ast)$
(Proposition~\ref{prop:B2-saga-traj} with $\muEff=\lmin(H_\ast)$, the value reached once the error is aligned). For
$\alpha\in\{0,1\}$ the column $9n\lmin(H_\ast)/(16R^2)$ of the table is below $1$ for the two runs that align, equal to
$0.977$ for the partially aligned run, and between $1.84$ and $9.05$ for the three runs that stay spread; with the global
$\mu$ in place of $\lmin(H_\ast)$, the value $9n/(16\kappa)=0.56$ at $\kappa=10^3$ would wrongly predict alignment for
$\alpha=0$. For $\alpha=2$ the column is below $1$ for $\kappa\ge10^3$, but the near-degeneracy of the bottom of the
spectrum makes alignment with a single eigenvector slow for both algorithms. One recorded quantity is consistent with
this dichotomy: the least-squares slope of $\log\norm{\theta_t-\etastar}^2$ against the number of passes, over the second
half of each run, is $-0.495$ and $-0.249$ in the two runs that align, i.e.\ the rate $2n\gamma\lmin(H_\ast)=0.493$ and
$0.249$ per pass of the slowest mode of the mean dynamics (to within $0.5\%$), whereas in the four runs with
$9n\lmin(H_\ast)/(16R^2)>1$ it lies between $-1.22$ and $-1.10$: a decay slower than the values $1.64$ to $8.05$ of
$2n\gamma\lmin(H_\ast)$ in these runs and close to the refresh rate of the memory (about $1$ per pass, a slot being
refreshed with probability $\frac1n$ at each iteration). This explanation is nevertheless only a heuristic: it relies on
the quadratic model, and the noise spectrum $c_k(t)$ was not recorded in these runs, so that the mechanism is inferred,
not measured; recording $q_{\min}^\top C_tq_{\min}$ and $\tr C_t$ (one pass over the data per evaluation) would turn the
consistency check above into a measurement.
\end{remark}

\subsection{How \texorpdfstring{$\muEff$}{mu\_eff} and \texorpdfstring{$\LEff$}{L\_eff} enter the SAGA bound}\label{sec:B2-saga-bound}

We use the setting of Proposition~5.9 of \cite{bach2024learning} (SAGA \cite{defazio2014saga} with
uniform sampling): step $\gamma=1/(4R^2)$, $\Delta=3\gamma^2$, and the Lyapunov function
\[
V_t=\norm{\theta_t-\etastar}^2+\Delta\sum_{i=1}^n\norm{z_i^{(t)}-\fip(\etastar)}^2 .
\]
Proposition~\ref{prop:B3-traj} of the B3 note is the following trajectory-dependent bound, obtained there by keeping
the co-coercivity step \eqref{ineq:B3-I2} and the strong-convexity step \eqref{ineq:B3-I4} of the book's proof as
identities; we include the short derivation because it shows exactly where $\muEff$ and $\LEff$ enter.

\begin{proposition}[trajectory-dependent one-step contraction]\label{prop:B2-saga-traj}
Under Assumption~\ref{ass:B2}, for every $t\ge1$, conditionally on the past and on the event $\theta_{t-1}\neq\etastar$,
\begin{equation}\label{ineq:saga-mueff-Leff}
\Ei\big[V_t\big]\;\le\;\Big(1-\min\Big\{\frac1{3n},\,a(\theta_{t-1})\Big\}\Big)V_{t-1},
\qquad a(\theta)=\frac{\muEff(\theta)\,\big(8R^2-5\LEff(\theta)\big)}{16R^4},
\end{equation}
and, since $\LEff(\theta)\le R^2$ implies $a(\theta)\ge3\muEff(\theta)/(16R^2)$,
\begin{equation}\label{ineq:saga-mueff}
\Ei\big[V_t\big]\;\le\;\Big(1-\min\Big\{\frac1{3n},\,\frac{3\muEff(\theta_{t-1})}{16R^2}\Big\}\Big)V_{t-1} .
\end{equation}
On the event $\theta_{t-1}=\etastar$ one has $\Ei[V_t]\le(1-\frac1{3n})V_{t-1}$. The bound of
Proposition~5.9, $\Ei[V_t]\le(1-\min\{\frac1{3n},\frac{3\mu}{16R^2}\})V_{t-1}$, follows from
\eqref{ineq:saga-mueff} and $\muEff\ge\mu$.
\end{proposition}

\begin{proof}
Write $v=\theta_{t-1}-\etastar$, $e_i=z_i^{(t-1)}-\fip(\etastar)$, and $\omega_t$ for the SAGA direction, so that
$\theta_t-\etastar=v-\gamma\omega_t$ and $\Ei[\omega_t]=\Fp(\theta_{t-1})$. Step~1 of the book's proof gives,
using only the expansion of the square, unbiasedness, $\norm{a+b}^2\le2\norm a^2+2\norm b^2$ and
$\Ei\norm{Z-\Ei Z}^2\le\Ei\norm Z^2$,
\[
\Ei\norm{\theta_t-\etastar}^2\le\norm v^2-2\gamma\,v^\top\Fp(\theta_{t-1})
+2\gamma^2\,\frac1n\sum_i\norm{\fip(\theta_{t-1})-\fip(\etastar)}^2+\frac{2\gamma^2}n\sum_i\norm{e_i}^2 ,
\]
and Step~2 gives the \emph{identity}
$\Ei\sum_i\norm{z_i^{(t)}-\fip(\etastar)}^2=(1-\frac1n)\sum_i\norm{e_i}^2+\frac1n\sum_i\norm{\fip(\theta_{t-1})-\fip(\etastar)}^2$.
Instead of bounding $\frac1n\sum_i\norm{\fip(\theta_{t-1})-\fip(\etastar)}^2$ by $R^2v^\top\Fp(\theta_{t-1})$
(co-coercivity), we use its exact value $\LEff(\theta_{t-1})\,v^\top\Fp(\theta_{t-1})$ (Definition~\ref{def:B2-eff}).
Adding $\Delta$ times the second relation to the first,
\[
\Ei[V_t]\le\norm v^2-2\gamma\Big(1-\gamma\LEff-\frac{\LEff\Delta}{2\gamma}\Big)v^\top\Fp(\theta_{t-1})
+\Big(\frac{2\gamma^2}{n\Delta}+1-\frac1n\Big)\Delta\sum_i\norm{e_i}^2 ,
\]
with $\LEff=\LEff(\theta_{t-1})$. With $\Delta=3\gamma^2$ and $\gamma=1/(4R^2)$:
$1-\gamma\LEff-\frac32\gamma\LEff=1-\frac52\gamma\LEff=1-\frac{5\LEff}{8R^2}\ge\frac38>0$ (as $\LEff\le R^2$),
and $\frac{2\gamma^2}{n\Delta}+1-\frac1n=1-\frac1{3n}$. Now, instead of $v^\top\Fp(\theta_{t-1})\ge\mu\norm v^2$,
we use the exact value $v^\top\Fp(\theta_{t-1})=\muEff(\theta_{t-1})\norm v^2$, and the coefficient of $\norm v^2$
becomes $1-2\gamma(1-\frac52\gamma\LEff)\muEff=1-\frac{\muEff(8R^2-5\LEff)}{16R^4}=1-a(\theta_{t-1})$, which is
nonnegative because $a\le\muEff/(2R^2)\le L/(2R^2)\le\frac12$. Hence
$\Ei[V_t]\le(1-a)\norm v^2+(1-\frac1{3n})\Delta\sum_i\norm{e_i}^2\le(1-\min\{\frac1{3n},a\})V_{t-1}$. If $v=0$,
the gradient-difference term vanishes and the same computation gives the factor $1-\frac1{3n}$.
\end{proof}

\begin{remark}[what each constant can explain]\label{rem:B2-what-explains}
With the step size fixed at $\gamma=1/(4R^2)$, the gain coming from $\LEff<R^2$ is the factor
$(8R^2-5\LEff)/(3R^2)\in[1,8/3]$ on the parameter-term exponent: the co-coercivity slack can explain at most a
factor $8/3$ of the gap, however small $\LEff$ is (a smaller $\LEff$ would allow a larger step, but the
algorithm does not take it). In contrast, the gain from $\muEff>\mu$ is the factor $\muEff/\mu$, which can be as
large as $L/\mu$ (Lemma~\ref{lem:B2-mueff}) and is typically of the order described in
Sections~\ref{sec:B2-logistic}--\ref{sec:B2-direction}. The parameter term is the active one in the minimum as
long as $a(\theta_{t-1})<1/(3n)$, i.e.\ (with $\LEff=R^2$) $\muEff(\theta_{t-1})<16R^2/(9n)$; in the regime
$\kappa\gg n$ of hypothesis (A4) this holds for $\mu$ and for moderate values of $\muEff$, and the bound saturates
at $1-\frac1{3n}$ once $\muEff$ exceeds $16R^2/(9n)$: beyond that point no improvement of the curvature term can be
``explained'', and the residual gap belongs to the memory coefficient $1-\frac1{3n}$, which is set by the square-norm
inequality \eqref{ineq:B3-I1}, the inequality $\Ei\norm{Z-\Ei Z}^2\le\Ei\norm Z^2$ \eqref{ineq:B3-I3} and the choice of
$\Delta$, and to the min step \eqref{ineq:B3-I5} between the two rates, all instrumented in the B3 note. In its
experiments with $\kappa\le10^2$, where the memory term is active from the start, order~A of
Proposition~\ref{prop:B3-attribution} (which, as here, keeps the minimum until the end) assigns $0.56$--$0.61$ of the
logarithmic gap to (I5), $0.38$--$0.44$ to (I1) and almost nothing to (I3).
\end{remark}

\paragraph{Explained fraction of the theory--practice gap.}
Fix a window $W=\{t_0,\dots,t_1-1\}$ of $T=t_1-t_0$ iterations inside the linear-convergence regime (e.g.\ one
effective pass, $T=n$), and let \texttt{V[t]}, \texttt{mu\_eff[t]}, \texttt{L\_eff[t]} be the recorded values of
$V_t$, $\muEff(\theta_t)$, $\LEff(\theta_t)$ (computed with $\etastar$ from a high-precision solver). Define the
per-iteration exponents (written $\nu$, to distinguish them from the tightness ratios $r_{\mathrm{sc}}$ and
$r_{\mathrm{coco}}$ recalled at the end of this subsection)
\begin{align}
\nu_{\mathrm{obs}}&=\frac1T\log\frac{V_{t_0}}{V_{t_1}}, \label{eq:B2-robs}\\
\nu_{\mathrm{th}}&=-\log\Big(1-\min\Big\{\frac1{3n},\frac{3\mu}{16R^2}\Big\}\Big), \label{eq:B2-rth}\\
\nu_{\mu}&=-\frac1T\sum_{t\in W}\log\Big(1-\min\Big\{\frac1{3n},\frac{3\muEff(\theta_t)}{16R^2}\Big\}\Big), \label{eq:B2-rmu}\\
\nu_{\mu L}&=-\frac1T\sum_{t\in W}\log\Big(1-\min\Big\{\frac1{3n},\frac{\muEff(\theta_t)\,(8R^2-5\LEff(\theta_t))}{16R^4}\Big\}\Big), \label{eq:B2-rmuL}
\end{align}
and the explained fractions
\begin{equation}\label{eq:B2-explained}
\phi_\mu=\frac{\nu_\mu-\nu_{\mathrm{th}}}{\nu_{\mathrm{obs}}-\nu_{\mathrm{th}}},\qquad
\phi_L=\frac{\nu_{\mu L}-\nu_\mu}{\nu_{\mathrm{obs}}-\nu_{\mathrm{th}}},\qquad
\phi_{\mathrm{rest}}=1-\phi_\mu-\phi_L=\frac{\nu_{\mathrm{obs}}-\nu_{\mu L}}{\nu_{\mathrm{obs}}-\nu_{\mathrm{th}}} .
\end{equation}
By Lemmas~\ref{lem:B2-mueff} and~\ref{lem:B2-Leff}, $a(\theta)\ge3\muEff(\theta)/(16R^2)\ge3\mu/(16R^2)$ and the monotonicity of
$x\mapsto-\log(1-\min\{\frac1{3n},x\})$, one has $\nu_{\mathrm{th}}\le \nu_\mu\le \nu_{\mu L}$, so
$\phi_\mu,\phi_L\ge0$ whenever $\nu_{\mathrm{obs}}>\nu_{\mathrm{th}}$ (the fractions are meaningful only when there is a gap
to explain). The quantity $e^{-T\nu_{\mu L}}=\prod_{t\in W}(1-\min\{\frac1{3n},a(\theta_t)\})$ is the product of the
one-step factors of Proposition~\ref{prop:B2-saga-traj} along the realized trajectory. It does not bound
$\E[V_{t_1}\mid\mathcal F_{t_0}]/V_{t_0}$ in general, because the factors are random, but the comparison holds in
expectation. Let $c_s=1-\min\{\frac1{3n},a(\theta_s)\}$, and $c_s=1-\frac1{3n}$ if $\theta_s=\etastar$; each $c_s$ is
$\mathcal F_s$-measurable ($\mathcal F_s=\sigma(i(1),\dots,i(s))$ as in Lemma~\ref{lem:B2-saga-quadratic}), and
Proposition~\ref{prop:B2-saga-traj} gives $\E[V_{s+1}\mid\mathcal F_s]\le c_sV_s$, so that
$Y_t=V_t/\prod_{s=t_0}^{t-1}c_s$, $t\ge t_0$, is a nonnegative supermartingale. Hence
$\E[e^{T(\nu_{\mu L}-\nu_{\mathrm{obs}})}\mid\mathcal F_{t_0}]=\E[Y_{t_1}\mid\mathcal F_{t_0}]/V_{t_0}\le1$ and, by
Jensen's inequality, $\E[\nu_{\mathrm{obs}}-\nu_{\mu L}\mid\mathcal F_{t_0}]\ge0$: the numerator of $\phi_{\mathrm{rest}}$
is nonnegative in conditional expectation, and for a single realization only up to fluctuations (average over seeds,
or use windows of several passes). Thus $\phi_\mu$ is the fraction of the gap (in exponent) explained by the
trajectory-dependent curvature alone, $\phi_L$ the additional fraction explained by the effective
co-coercivity constant (at most a factor $8/3$ by Remark~\ref{rem:B2-what-explains}), and $\phi_{\mathrm{rest}}$ the
fraction left to the other inequalities of the proof (square-norm inequality, variance versus second moment,
the minimum of the two rates, the choice of $\Delta$), each of which the B3 note measures separately. The same
decomposition can be reported instantaneously through the gain factors
\[
g_\mu(t)=\frac{\min\{\frac1{3n},\frac{3\muEff(\theta_t)}{16R^2}\}}{\min\{\frac1{3n},\frac{3\mu}{16R^2}\}}\ \ge1,\qquad
g_L(t)=\frac{\min\{\frac1{3n},a(\theta_t)\}}{\min\{\frac1{3n},\frac{3\muEff(\theta_t)}{16R^2}\}}\in\Big[1,\frac83\Big],
\]
whose product is the ratio between the guaranteed per-step exponent along the trajectory and the one of
Proposition~5.9. Finally, the per-step tightness ratios of the two inequalities themselves are, in the notation of
Definition~\ref{def:B3-ratios} of the B3 note,
\begin{align*}
r_{\mathrm{sc}}(t)&=\frac{\mu\norm{\theta_t-\etastar}^2}{(\theta_t-\etastar)^\top\Fp(\theta_t)}=\frac{\mu}{\muEff(\theta_t)}\in\Big[\frac\mu L,1\Big],\\
r_{\mathrm{coco}}(t)&=\frac{\frac1n\sum_i\norm{\fip(\theta_t)-\fip(\etastar)}^2}{R^2(\theta_t-\etastar)^\top\Fp(\theta_t)}=\frac{\LEff(\theta_t)}{R^2}\in\Big[\frac\mu{R^2},1\Big],
\end{align*}
for \eqref{ineq:mueff-lower} and \eqref{ineq:Leff-cocoercive} respectively.

\subsection{Related work}\label{sec:B2-related}

\emph{Restricted secant inequality.} Since the minimizer is unique, $\muEff(\theta)$ is the constant, at the point
$\theta$, of the restricted secant inequality of \cite{zhang2013gradient} (Definition~2 there:
$\inner{\Fp(x)}{x-x_{\mathrm{prj}}}\ge\nu\norm{x-x_{\mathrm{prj}}}^2$, with $x_{\mathrm{prj}}$ the projection of $x$ on the
solution set), whose constant ``can be viewed as a lower bound of the average curvature'' of the function between $x$
and $x_{\mathrm{prj}}$; Lemma~\ref{lem:B2-integral} makes this exact. The best global constant of this inequality is
$\inf_{\theta\neq\etastar}\muEff(\theta)$, which equals $\lambda$ for logistic regression: along $\theta=\etastar+t\hat d$,
$\hat d$ a unit vector, \eqref{eq:B2-logistic-mueff} gives $\muEff(\theta)-\lambda=\frac1n\sum_iw_i(\theta)(x_i^\top\hat d)^2$,
where $w_i(\theta)\le1/\abs{a_i-a_i^\ast}=1/(t\abs{x_i^\top\hat d})$ by \eqref{eq:B2-weights-closed} when
$x_i^\top\hat d\neq0$, so that $\muEff(\etastar+t\hat d)\to\lambda$ as $t\to\infty$.

\emph{Local strong convexity and self-concordance.} The observation that the curvature actually experienced by an
algorithm is the local one is classical; the book itself notes that ``locally around the global optimum, the
strong convexity constant converges to the minimal eigenvalue of the Hessian at $\etastar$, which can be
significantly larger than $\mu$ (the global constant)'' (\cite{bach2024learning}, Section~5.2.3, paragraph
``Adaptivity''). For logistic regression, \cite{bach2014adaptivity} proves that averaged SGD with constant
step size proportional to $1/(R^2\sqrt N)$ has an $O(1/\sqrt N)$ rate which ``improves to $O(R^2/(\mu N))$ where
$\mu$ is the lowest eigenvalue of the Hessian at the global optimum'' (when this eigenvalue is greater than
$R^2/\sqrt N$), without knowledge of $\mu$
(adaptivity to unknown local strong convexity); the tool is the generalized self-concordance of the logistic
loss, Assumption (A2) there: along every line, $\abs{\varphi'''(t)}\le R\norm{\theta_1-\theta_2}\varphi''(t)$, a
notion introduced in \cite{bach2010self}, whose Proposition~1, eq.~(6), gives the Hessian comparison
$e^{-R\norm v}F''(w)\preccurlyeq F''(w+v)\preccurlyeq e^{R\norm v}F''(w)$. Lemma~\ref{lem:B2-selfconc} and
Corollary~\ref{cor:B2-local} are the per-sample, directional version of this comparison, and $\muEff(\theta_t)$ is
the trajectory-wise counterpart of the local constant $\lmin(H_\ast)$: by Proposition~\ref{prop:B2-limits} it
converges to the Rayleigh quotient $u_t^\top H_\ast u_t\ge\lmin(H_\ast)$, with equality only when the error aligns
with the flattest direction, which Section~\ref{sec:B2-direction} predicts for gradient descent but not for
SAGA (the experiments confirm the latter only for $\kappa$ at most of the order of $n$, and not in every such run,
Remark~\ref{rem:B2-P2-outcome}).

\emph{Expected smoothness.} The constant $\LEff$ is a trajectory-wise version of the \emph{expected
smoothness} constant of \cite{gower2021stochastic} (Assumption~3.1 there: there is $\mathcal L_1>0$ with
$\E_{\mathcal D}\norm{\nabla f_S(x)-\nabla f_S(x^\ast)}^2\le2\mathcal L_1(f(x)-f(x^\ast))$ for all $x$), which
governs the rates of SAGA and of its JacSketch generalizations; for minibatch sketches with $\tau$-nice sampling,
Theorem~2 there (Theorem~4.16 of the arXiv version 1805.02632v1) shows that Assumption~3.1 is satisfied with
$\mathcal L_1=L^{\mathcal G}_{\max}$ (eq.~(85) there), which for single-element uniform sampling ($\tau=1$) is $\max_iL_i$,
a valid constant rather than the smallest one; the same assumption is the basis of the general SGD analysis of
\cite{gower2019sgd} (Assumption~2.1: $\E_{\mathcal D}\norm{\nabla f_v(x)-\nabla f_v(x^\ast)}^2\le2\mathcal L(f(x)-f(x^\ast))$).
Defining the trajectory-wise expected smoothness
$\mathcal L_{\mathrm{ES}}(\theta)=\frac1n\sum_i\norm{\fip(\theta)-\fip(\etastar)}^2/\big(2(F(\theta)-F(\etastar))\big)$,
the convexity inequality $F(\theta)-F(\etastar)\le(\theta-\etastar)^\top\Fp(\theta)$ (eq.~(5.7) of the book) gives
$\mathcal L_{\mathrm{ES}}(\theta)\ge\LEff(\theta)/2$, with $\mathcal L_{\mathrm{ES}}=\LEff$ for quadratics (where
$2(F(\theta)-F(\etastar))=(\theta-\etastar)^\top\Fp(\theta)$). We work with $\LEff$ because the proof of
Proposition~5.9 is written with the inner product $(\theta-\etastar)^\top\Fp(\theta)$ rather than with function values;
both constants measure how much of the worst-case $R^2=\max_iL_i$ is used along the trajectory, and both can
be recorded at the same cost.
